\documentclass[11pt]{article}
\usepackage[T1]{fontenc}
\usepackage[utf8]{inputenc}
\usepackage{lmodern,microtype,amsmath,amssymb,amsthm,mathtools,booktabs,longtable,listings,needspace}
\usepackage[margin=1in]{geometry}
\usepackage[authoryear,round]{natbib}
\usepackage[hidelinks]{hyperref}
\setlength{\emergencystretch}{2em}
\lstset{basicstyle=\ttfamily\footnotesize,breaklines=true,columns=fullflexible,keepspaces=true,showstringspaces=false}
\newtheorem{theorem}{Theorem}[section]
\newtheorem{proposition}[theorem]{Proposition}
\newtheorem{lemma}[theorem]{Lemma}
\theoremstyle{remark}\newtheorem{remark}[theorem]{Remark}
\newcommand{\Z}{\mathbb Z}\newcommand{\Q}{\mathbb Q}\newcommand{\F}{\mathbb F}
\newcommand{\Oa}{\mathcal O}\newcommand{\Cl}{\operatorname{Cl}}
\newcommand{\Gal}{\operatorname{Gal}}\newcommand{\len}{\operatorname{length}}
\newcommand{\coker}{\operatorname{coker}}\newcommand{\lc}{\widetilde{\Cl}_3}
\newcommand{\ind}{\mathbf1}
\newcommand{\workbenchurl}{https://github.com/KokunoYumeto/zeta-and-l-functions-workbench}
\newcommand{\recordurl}{https://github.com/KokunoYumeto/mathematics-workbenches}
\title{The $3$-primary $K$-groups of $\Q(\omega,\sqrt[3]{269})$ in every degree\\[3pt]
\large After Schneider and Jaulent--Michel}
\author{}
\date{24 November 2026}
\begin{document}\maketitle
\begin{abstract}
Schneider's description of wild kernels by cyclotomic coinvariants, developed in logarithmic form by Jaulent and Michel, reduces even algebraic $K$-groups to arithmetic in a cyclotomic tower.
For $E=\Q(\omega,\sqrt[3]{269})$, where $\omega$ is a primitive cube root of unity, we determine
\[
 K_{2e}(\Oa_E)_{(3)}\simeq\Z_3[\omega]/(\omega-1)^{5+2v_3(e)+\ind_{e\equiv2\ (3)}}
 \qquad(e\geq1)
\]
from finite logarithmic class-group certificates in degrees six, nine, and eighteen.
A ray-class genus calculation determines the residual cyclotomic module, a short specialization lemma uses the first-layer logarithmic bound, and a labelled action on a cyclic group of order nine determines the exceptional residue class.
We also prove the precise length-one correction between the character-valued strict group over $\Q(\omega)$ and the ordinary strict group over $E$, including vanishing of the former at weight zero.
A finite theta estimate proves positivity of the associated completed cubic Hecke function on the real axis.
An appendix also computes the faithful seven-primary
cyclotomic class module and every continuous cyclotomic strict specialization.
A further appendix determines the one-prime ramified module: it has a
linear characteristic factor whose values are unbounded along positive
integer weights, while the matched strict lengths remain one.
The text and programs were produced by GPT-6 Astra in ChatGPT Pro mode and checked by separate instances of the same model as recorded in the workbench; no human verification is claimed.
\end{abstract}
\tableofcontents

\section{Introduction and main result}
This paper determines every positive even three-primary $K$-group of one explicit nonabelian number field and explains the arithmetic responsible for its variation with the degree.
Schneider's cyclotomic description of wild kernels \citep{Schneider1979} and its logarithmic development by \citet{JaulentMichel2006} provide the starting point.
The comparison with algebraic $K$-theory and the global-to-local sequence are standard theorems; we use the formulations of \citet{Weibel2013} and \citet{KolsterMovahhedi2000}.
The logarithmic algorithms of \citet{BelabasJaulent2016} supply the finite arithmetic calculations.
The additional argument here proves that a small amount of finite information controls the entire cyclotomic module, determines the action with its correct orientation, and thereby gives an explicit answer in every degree.

The field has no three-primary ordinary ideal classes, but this does not force its wild kernels to vanish.
The strict character group over the quadratic subfield is also distinct from the ordinary strict group over the full field; their lengths differ by one.
Both distinctions are necessary for the formula below.
The $K$-theoretic assertions concern the coefficient prime three and the single field $E$. A separate analytic theorem adapts the finite partial-sum method of \citet{LouboutinMunsch2013} to the associated cubic Hecke function. It proves positivity on the positive real axis; the nonreal zeros remain outside its scope.

Recent work over the same Eisenstein field studies the associated analytic questions.
\citet[Theorem~1.1, arXiv version~3, p.~3]{DunnRadziwill2024} prove Patterson's prime Gauss-sum asymptotic under Hecke GRH.
\citet[Theorem~1.1, p.~1]{OpenAIFirst2026} give an unconditional first-moment asymptotic, and \citet[Theorem~1.1, p.~4]{OpenAIQuasi2026} give a zero-free half-plane $\Re s>7/8$ for finite-order Hecke functions over $K$.
Theorem~\ref{thm:cubic-real-positive} below proves real nonvanishing for the specific cubic factor by an explicit positive kernel.
The every-degree $K$-group formula rests on logarithmic descent and uses none of these prime-sum or zero-free estimates.

\subsection{Conventions}
Put
\[
 K=\Q(\omega),\quad \alpha^3=269,\quad E=K(\alpha),\quad
 O=\Z_3[\omega],\quad L=\Q_3(\omega),\quad D=L/O.
\]
Choose $\sigma\in\Gal(E/K)$ by $\sigma\alpha=\omega\alpha$.
When the cubic norm vanishes, the induced $O$-action means $\sigma\mapsto\omega$.
Normalize $v_3(3)=1$, $v_\lambda(\lambda)=1$, and put
\[
 \lambda=\omega-1,\quad \lambda^2=-3\omega,\quad
 \delta=\omega-\omega^2,\quad\delta^2=-3.
\]
The elements $\delta$ and $\lambda$ differ by a unit.
For finite $O$-modules, $O$-length and $\Z_3$-length both equal the base-three logarithm of the order.
The notation $A_{(3)}$ denotes the entire three-primary subgroup.

Write $\varepsilon$ for the three-adic cyclotomic character, with $\Z_3(1)$ carrying action $\varepsilon$.
Set $\chi(g)=g\alpha/\alpha$ and $D_e=D(\chi\varepsilon^{-e})$.
Let $S$ consist of the primes of $K$ above $3$ and $269$.
The character-valued strict group is
\begin{equation}
 \mathcal S_e=\ker\left(H^1(G_{K,S},D_e)\longrightarrow
 H^1(K_3,D_e)\oplus H^1(I_{269},D_e)\right),
 \label{eq:charS}
\end{equation}
where $K_3=\Q_3(\omega)$ and $I_{269}$ is absolute inertia.
For $e\geq1$ the ordinary strict group over $E$ is
\begin{equation}
 U_e=\ker\left(H^1(G_{E,S_3},\Q_3/\Z_3(-e))
 \longrightarrow\bigoplus_{w\mid3}H^1(E_w,\Q_3/\Z_3(-e))\right).
 \label{eq:ordinaryS}
\end{equation}
All cohomology is continuous. These are different groups, over different base fields and coefficient rings.

\begin{theorem}[the every-degree formula]
\label{thm:main}
For every $e\geq1$, put
\[
 m_e=4+\ind_{e\equiv2\ (3)},\quad a_e=3+\ind_{e\equiv2\ (3)},\quad
 \kappa_e=5+2v_3(e)+\ind_{e\equiv2\ (3)}.
\]
Then
\[
 K_{2e}(\Oa_E)_{(3)}\simeq O/\lambda^{\kappa_e},\qquad
 U_e^\vee\simeq O/\lambda^{a_e},\qquad
 \mathcal S_e\simeq O/\lambda^{m_e}.
\]
The first two actions are induced by $\sigma\mapsto\omega$, with the dual action understood in the second.
At weight zero and for ordinary classes,
\[
 \Cl(E)_{(3)}=0,\qquad\mathcal S_0=0.
\]
For $N\geq0$,
\[
 \sum_{e=1}^N\len_O\mathcal S_e
 =4N+\left\lfloor\frac{N+1}{3}\right\rfloor.
\]
\end{theorem}
The standard strict group over $E$ has length $3$ or $4$; the length $4$ or $5$ belongs to the character group over $K$.
The degree-zero logarithmic class group at weight zero, used later, is a third object. It is not $\mathcal S_0$.
The four finite arithmetic assertions used in the proof are proved in Section~\ref{sec:certificates}, with complete computational specifications in the appendices.

\begin{theorem}[the associated cubic Hecke function]
\label{thm:cubic-real-positive}
Let $F=\Q(\alpha)$, $Q=217083$, $a=2\pi/\sqrt Q$, and
\[
 L_F^{\mathrm{rel}}(s)=\frac{\zeta_F(s)}{\zeta(s)}
   =\sum_{n\ge1}A(n)n^{-s}\quad(\Re s>1).
\]
The completed entire function
\[
 \Lambda_F^{\mathrm{rel}}(s)=2a^{-s}\Gamma(s)L_F^{\mathrm{rel}}(s)
\]
is strictly positive for every real $s$, and $L_F^{\mathrm{rel}}(s)>0$ for every $s>0$.
Moreover
\[
 \sum_{n\ge1}A(n)e^{-ant}
       \ge\frac{999999}{1000000}e^{-at}>0\qquad(t\ge1).
\]
The statement concerns real values and does not assert a Riemann hypothesis.
\end{theorem}

\subsection{A related exact theorem at seven}

The finite arithmetic at seven gives a separate stability phenomenon for
the same cubic extension. Its complete proof and exact witnesses are
included in Appendix~\ref{seven:section}.
We state the result here with its local conditions and coefficient
normalization; the three-primary $K$-group formula does not depend on it.

\begin{theorem}[seven-primary stability and strict specialization]
\label{thm:seven-stability}
Let $Q_n^{(7)}$ be the degree-$7^n$ layer of the cyclotomic
$\Z_7$-extension of $\Q$, and put
$F_n^{(7)}=\Q(\alpha)Q_n^{(7)}$ and
$E_n^{(7)}=EQ_n^{(7)}$.
Let $\psi$ be either faithful character of $\Delta=\Gal(E/K)$,
with values in $\mu_3\subset\Z_7^\times$.
Choose $\gamma_7\in\Gamma_7=\Gal(E_\infty^{(7)}/E)$ with
$\langle\varepsilon_7\rangle(\gamma_7)=8$, and put $T_7=\gamma_7-1$.
Then, for every $n\ge0$,
\[
 \Cl(F_n^{(7)})_{(7)}\simeq\Z/7\Z,\qquad
 \left(\varprojlim_n\Cl(E_n^{(7)})_{(7)}\right)^\psi
 \simeq\Z_7[[T_7]]/(7,T_7).
\]
The norms on the first family are isomorphisms and its cyclotomic action
is trivial.

Let $V_7/\Q_7$ be finite, with integer ring $\Oa_7$, and let
$\kappa:\Gamma_7\to\Oa_7^\times$ be continuous.
Define the strict group for
$D_{\psi,\kappa}=(V_7/\Oa_7)(\psi\kappa)$ by zero restriction to both full
decomposition groups above seven and to inertia at every other finite
place of $K$. Then
\[
 \operatorname{Sel}_{\mathrm{str},7}(K,D_{\psi,\kappa})
 \simeq\Oa_7/(7,\kappa(\gamma_7)-1).
\]
In particular, for every $e\in\Z_7$, including zero, the principal
character $\kappa=\langle\varepsilon_7\rangle^{-e}$ gives
$\operatorname{Sel}_{\mathrm{str},7}\simeq\Oa_7/7\Oa_7$ and
\[
 \frac{\log\#\operatorname{Sel}_{\mathrm{str},7}}
      {[V_7:\Q_7]}=\log7 .
\]
\end{theorem}



The coefficient normalization matters when $V_7/\Q_7$ is ramified:
$\Oa_7/7\Oa_7$ has $\Oa_7$-length $e(V_7/\Q_7)$, while its displayed
normalized logarithmic size remains $\log7$.
If $\kappa(\gamma_7)=\zeta_{7^r}\exp(-e\log8)$ with $r\ge1$,
$e\in\Oa_7$, and $\zeta_{7^r}$ primitive, the same theorem gives
$\Oa_7/(\zeta_{7^r}-1)$ and normalized size
$\log7/(6\cdot7^{r-1})$.
These conclusions do not assert vanishing of the trivial
cubic-character component or a bound for other arithmetic families.

\begin{theorem}[the one-prime module and its varying specializations]
\label{thm:seven-one-prime}
Retain the notation of Theorem~\ref{thm:seven-stability}, and put
$\Lambda_7=\Z_7[[T_7]]$. Fix either prime $\mathfrak p\mid7$ of $K$.
Let $U_{\mathfrak p,\infty}$ be the norm inverse limit of principal units
in the completions of $E_n^{(7)}$ above $\mathfrak p$, and let
$\overline{\mathcal E}_\infty$ be the inverse limit of the closures of
localized global units. Put
$V_{\mathfrak p,\psi}=(U_{\mathfrak p,\infty}/
\overline{\mathcal E}_\infty)^\psi$.
Let $X_{\mathfrak p,\infty}$ be the Galois group of the maximal abelian
pro-seven extension of $E_\infty^{(7)}$ unramified outside the primes
over $\mathfrak p$. There exists $a_{\mathfrak p,\psi}\in7\Z_7^\times$
such that
\[
 V_{\mathfrak p,\psi}\simeq X_{\mathfrak p,\infty}^\psi
 \simeq\Lambda_7/(T_7-a_{\mathfrak p,\psi}).
\]
In the natural class-field sequence
\[
 0\longrightarrow V_{\mathfrak p,\psi}
 \longrightarrow X_{\mathfrak p,\infty}^\psi
 \longrightarrow\Lambda_7/(7,T_7)\longrightarrow0,
\]
the image of the first map is $7X_{\mathfrak p,\infty}^\psi$.
Set $e_{\mathfrak p,\psi}^*=-\log(1+a_{\mathfrak p,\psi})/\log8$.
Then $e_{\mathfrak p,\psi}^*\in\Z_7^\times$, and every characteristic
ideal generator $f_{\mathfrak p,\psi}$ satisfies
\[
 v_7\bigl(f_{\mathfrak p,\psi}(8^{-e}-1)\bigr)
 =1+v_7(e-e_{\mathfrak p,\psi}^*)\qquad(e\in\Z_7),
\]
with the value infinity at $e=e_{\mathfrak p,\psi}^*$.
These valuations are unbounded for positive integers $e$, whereas
the strict groups with coefficients $\Q_7/\Z_7$ in
Theorem~\ref{thm:seven-stability} have length one at every such weight.
\end{theorem}

The full proof, including the exact first-layer local-unit matrices,
is in Appendix~\ref{seven:universal-norms}. It uses Rubin's class-field
sequence and his absence of finite submodules for the one-prime module
\citep[equation~(5), p.~36; Theorem~5.3(v), pp.~40--41]{Rubin}.
No identification with a separately normalized Katz function is needed
or asserted.

\section{Proof overview}
The proof begins with local arithmetic and the ambiguous class formula, which show that the prime above three splits completely in the cubic extension and that its three-primary ordinary class group is zero.
Poitou--Tate duality then isolates a local correction of length $2(1+v_3(e))$.
The residual dimension and the vanishing cubic norm force the entire even $K$-group to be cyclic over $\Z_3[\omega]$.
A separate restriction argument compares the two strict Selmer groups and proves their exact length-one difference.

The remaining problem is to determine the wild kernel in every cyclotomic twist.
There are nine primes above $269$ in the cyclotomic tower of $\Q(\omega)$.
Their residue-unit groups contribute nine generators to a ray module, and the roots of unity give one diagonal relation.
The condition that primes above three split adds no further relation: the negative part of the group of three-units already consists of roots of unity.
Passing to the cubic extension removes one further generator and gives the exact residual module $\F_3[T]/T^7$, where $T=\gamma-1$.
This field-specific genus information is essential to the argument.

Complex conjugation separates a cyclic presentation of the cyclotomic module into two ideals over $\Z_3[[T]]$.
The positive component of the first logarithmic layer is cyclic of order nine, with $\gamma$ acting by four.
The order bound for the entire layer then forces the other component to have exponent nine.
Lifting this exponent relation to the power-series ring and using the genus information shows that its value at every $T\in3\Z_3$ generates the ideal $(9)$.
The positive component contributes order nine exactly when $T\equiv3\pmod9$, and order three otherwise.
The twist has $T=4^{-e}-1$, so the exceptional condition is $e\equiv2\pmod3$.
Adding the local correction and the character-descent correction gives the two formulas.
The same ideal argument also proves finiteness of the full cyclotomic module.

The analytic assertion has a separate proof. Two exact coefficient constructions identify the cubic Hecke series, and finite positive partial sums control its theta kernel after a uniform geometric tail estimate. The functional equation turns this kernel into a positive Mellin integral on the real axis. The same integral supplies a rigorously bounded value at one, used with exact ideal relations to determine the ordinary class groups.
The seven-primary theorem uses Appendix~\ref{seven:section} and its two finite
class-group certificates and a separate one-prime norm-control argument;
it is independent of the three-adic logarithmic specialization.
For the one-prime module, class capitulation measures a persistent index
seven in norms of global units. Local norms have exact control on the
faithful character. A printed first-layer local determinant then forces
the unit quotient to have a single linear characteristic factor.
Rubin's theorem excludes a finite submodule in the ramified Galois
module, so the unit quotient embeds as multiplication by seven.
Evaluating that linear factor gives the moving seven-adic center in
Theorem~\ref{thm:seven-one-prime}.

\section{Ramification, ordinary classes, and cyclic \texorpdfstring{$K$}{K}-groups}
This section proves the local splitting, the ordinary class-group vanishing, and the cyclic structure needed to reduce the problem to the wild kernel.
\begin{lemma}[local arithmetic]
\label{lem:local}
The cyclic cubic extension $E/K$ is totally ramified at $(269)$ and unramified at every other finite prime.
The prime above three splits completely, each completion $E_w$ above three is $K_3$, and
$E\cap\Q(\mu_{3^\infty})=K$.
\end{lemma}
\begin{proof}
The prime $269\equiv2\pmod3$ is inert in $K$ and $X^3-269$ is Eisenstein there.
Its discriminant is supported at three and $269$.
Hensel's lemma applied to $h(t)=t-3t^2+3t^3-30$ gives a root in $\Z_3$, since $h(0)\equiv0$ and $h'\equiv1\pmod3$.
Then $(-1+3t)^3=269$, so all three roots are in $K_3$.
Thus three splits and introduces no ramification.
A nontrivial cyclotomic intersection would contain all of $E$, because $[E:K]=3$, contradicting ramification at $269$.
\end{proof}

\begin{proposition}[ordinary class-group vanishing]
\label{prop:classes}
One has $\Cl(E)_{(3)}=0$.
\end{proposition}
\begin{proof}
The norm-Euclidean ring $\Z[\omega]$ has class number one.
Chevalley's ambiguous class formula in the totally imaginary case gives
\[
 \#\Cl(E)^{\langle\sigma\rangle}
 =\frac{h_K\prod_v e_v}{3[\Oa_K^\times:\Oa_K^\times\cap N_{E/K}E^\times]}
 =\frac1{[\Oa_K^\times:\Oa_K^\times\cap N_{E/K}E^\times]}.
\]
The ordinary-class formula with this norm index is Theorem~1.1 of \citet{Yu2014}.
The left side is a positive integer, so it is one.
If a nonzero finite three-group has an order-three automorphism, that automorphism has a nonzero fixed vector on its three-torsion: $(\sigma-1)^3=0$ in characteristic three.
Applying this to $\Cl(E)_{(3)}$ proves the claim.
\end{proof}

\begin{proposition}[cyclicity and the local correction]
\label{prop:structure}
For $e\geq1$ write $H_e=K_{2e}(\Oa_E)_{(3)}$, $\kappa_e=\log_3\#H_e$, and $k_e=1+v_3(e)$.
Then $H_e$ is an $O$-module and
\begin{equation}
 H_e\simeq O/\lambda^{\kappa_e},\qquad
 U_e^\vee\simeq O/\lambda^{\kappa_e-2k_e}.
 \label{eq:structure}
\end{equation}
In particular $\kappa_e\geq2k_e$.
\end{proposition}
\begin{proof}
The integral \'etale comparison identifies $H_e$ with $H^2(G_{E,S_3},\Z_3(e+1))$.
We use \citet[Chapter VI, Theorem 8.2, Corollary 8.3, and Example 8.3.1]{Weibel2013} for the comparison and its residual form.
Poitou--Tate duality gives
\begin{equation}
\begin{split}
0\longrightarrow U_e^\vee\longrightarrow H_e\longrightarrow
 \bigoplus_{w\mid3}H^0(E_w,\Q_3/\Z_3(-e))^\vee\\
 \longrightarrow H^0(E,\Q_3/\Z_3(-e))^\vee\longrightarrow0.
\end{split}\label{eq:PT}
\end{equation}
This wild-kernel sequence is also recorded in \citet[Section 1]{KolsterMovahhedi2000}.
All four global and local invariant groups have order $3^{k_e}$, because their cyclotomic images are $1+3\Z_3$ and
\[
 \min_{u\in1+3\Z_3}v_3(u^e-1)=1+v_3(e),
\]
attained at $u=4$.
Taking orders yields
\begin{equation}
 \len_{\Z_3}U_e=\kappa_e-2k_e.\label{eq:correction}
\end{equation}
For $K$, the residual comparison gives $K_{2e}(\Oa_K)/3=0$: the localized class group is trivial and there is one prime above three.
Its three-primary group is therefore zero.
Transfer to $K$ followed by restriction makes $1+\sigma+\sigma^2$ vanish on $H_e$.
One may first use the transfer identity on fraction fields; the localization sequence injects the even positive integer-ring $K$-group into its field $K$-group, because the even positive $K$-groups of finite residue fields vanish.
This justifies the norm identity even at the ramified prime.
It induces the asserted $O$-action.

The residual comparison, Proposition~\ref{prop:classes}, and the three primes above three give $\dim_{\F_3}H_e/3H_e=2$.
A lift of $\sigma$ fixing all cyclotomic roots cyclically permutes the three local terms in \eqref{eq:PT}; the last map is summation.
Its kernel is the augmentation module over $\Z/3^{k_e}$, hence is $O/3^{k_e}O=O/\lambda^{2k_e}$.
Thus $H_e$ surjects onto this module.
Write $H_e=\bigoplus_iO/\lambda^{n_i}$.
The residual dimension says $\sum_i\min(2,n_i)=2$.
The surjection forces one $n_i\geq2k_e\geq2$, accounting for the whole dimension; no other summand exists.
The kernel in \eqref{eq:PT} is then cyclic of length \eqref{eq:correction}.
\end{proof}

\section{Strict character descent and weight zero}
This section compares the two strict groups with their exact coefficient conventions and proves the additional length-one term.
\begin{lemma}[the residual strict group]
\label{lem:residual}
One has $\mathcal S_0=0$.
For every nonzero integer $e$, $\mathcal S_e[\lambda]\simeq\F_3$, represented by the Kummer class of $\omega$.
For $e\geq1$, $\mathcal S_e$ is finite cyclic over $O$.
\end{lemma}
\begin{proof}
Kummer theory and the class number of $K$ give
\[
 H^1(G_{K,S},\F_3)=\F_3[\omega]\oplus\F_3[\lambda]\oplus\F_3[269].
\]
Restriction at three kills $[269]$ by Lemma~\ref{lem:local}.
The other two classes are independent: $\lambda$ has valuation one, and $\omega$ is a noncube unit, since its cube root generates $\Q_3(\zeta_9)/K_3$.
At $269$, the pro-three tame inertia generator acts on $D_e$ by $\omega$.
Multiplication by $\lambda$ is surjective on $D$, so every cocycle there is a coboundary.
Prime-to-three inertia contributes no positive cohomology, whence
\begin{equation}H^1(I_{269},D_e)=0.\label{eq:inertia}
\end{equation}
The exact coefficient sequence $0\to\F_3\to D_e\xrightarrow{\lambda}D_e\to0$ gives
\[
 H^1(H,D_e)[\lambda]=H^1(H,\F_3)/\partial(D_e^H/\lambda D_e^H).
\]
Globally, ramified inertia forces $D_e^{G_{K,S}}=D_e[\lambda]$.
Lifting $1/\lambda$ to $1/\lambda^2$ identifies its connecting line with $\F_3[269]$; the cyclotomic factor is one modulo $\lambda^2$.
Locally at three the connecting line is zero for $e=0$, since the invariant module is divisible.
For $e\ne0$, put $a=1+v_3(e)$; the local invariant module is $D[\lambda^{2a}]$.
Lifting its generator through multiplication by $\lambda$ gives, up to a unit, the character
\[
 \frac{\varepsilon(g)^{-e}-1}{3^a}
 \equiv-\frac{e}{3^{a-1}}\frac{\varepsilon(g)-1}{3}\pmod3.
\]
The first scalar is a three-adic unit, even when three divides $e$.
Thus the connecting character has Kummer line $\F_3[\omega]$.
Naturality therefore makes the strict socle zero at weight zero and the $\omega$-line otherwise.
Every continuous cocycle has finite image; hence the cohomology is primary torsion, and a zero socle implies the whole group is zero.

For $e\geq1$, restriction to $E$ has image in the finite group $U_e\otimes_{\Z_3}O$, by \eqref{eq:PT}.
Its kernel lies in $H^1(\Gal(E/K),D_e^{G_E})$, which is finite because $e>0$ makes the cyclotomic invariants finite.
Thus $\mathcal S_e$ is finite.
A finite module over a discrete valuation ring with one-dimensional socle is cyclic.
\end{proof}

\begin{proposition}[exact length-one correction]
\label{prop:descent}
For $e\geq1$, restriction gives
\begin{equation}
0\longrightarrow\F_3\longrightarrow\mathcal S_e\longrightarrow
 (U_e\otimes_{\Z_3}O(\chi))^{\langle\sigma\rangle}\longrightarrow0,
\label{eq:descent}
\end{equation}
whose target has order $\#U_e$.
Consequently
\begin{equation}
\len_O\mathcal S_e=1+\len_{\Z_3}U_e
 =\kappa_e-2(1+v_3(e))+1.\label{eq:exactlength}
\end{equation}
\end{proposition}
\begin{proof}
With $k=1+v_3(e)$, the invariant coefficient $N=D_e^{G_{E,S}}$ is $D_e[\lambda^{2k}]$, on which $C_3=\Gal(E/K)$ acts by $\omega$.
Its periodic resolution gives $H^1(C_3,N)=N/\lambda N=\F_3$ and $H^2(C_3,N)=N[\lambda]=\F_3$.
To prove surjectivity onto the strict invariant target, consider the map from
$1\to I_E\to I_K\to C_3\to1$ at $269$ to the global extension sequence.
The coefficient on $I_E$ is trivial, and
\[
 H^2(C_3,N)\longrightarrow H^2(C_3,D_e)=D_e[\lambda]
\]
is the inclusion of the same socle, hence an isomorphism.
A strict class restricts to zero on $I_E$; naturality sends its global transgression to zero in this local group.
The transgression is therefore zero.
Every invariant strict class lifts globally; its localization at three is zero because the completion extension is the identity, and its inertia condition follows from \eqref{eq:inertia}.
The inflated kernel is strict for the same reasons, proving \eqref{eq:descent}.

The norm $1+\sigma+\sigma^2$ kills $U_e$, by dualizing the corresponding norm statement in Proposition~\ref{prop:structure}.
On $P_0=U_e\otimes_{\Z_3}O$, let $q=\sigma\otimes\omega-1$.
The invariant target is $\ker q$, while
\[
\coker q=U_e\otimes_{\Z_3[C_3]/(1+\sigma+\sigma^2),\,\sigma\mapsto\omega^{-1}}O.
\]
The coefficient map is an isomorphism, conjugating the two cubic roots, so this cokernel has order $\#U_e$.
Kernel and cokernel of a finite-group endomorphism have equal orders.
Taking lengths in \eqref{eq:descent} and using \eqref{eq:correction} proves the result.
\end{proof}

\section{Cyclotomic descent and the ray genus module}
This section realizes the wild kernels as twisted coinvariants and computes the residual cyclotomic module by ray class theory.
Let
\[
 K_n=\Q(\mu_{3^{n+1}}),\quad E_n=EK_n,\quad
 A_n=\Cl(E_n)_{(3)},\quad A'_n=\Cl(\Oa_{E_n}[1/3])_{(3)}.
\]
Use norm transitions to define $M=\varprojlim A'_n$.
Equivalently, $M$ is the Galois group over $E_\infty$ of its maximal abelian pro-three extension unramified at every finite prime and split completely at primes above three.
Normalize $\gamma\in\Gamma=\Gal(E_\infty/E)$ by $\varepsilon(\gamma)=4$, and put
\[
 T=\gamma-1,\qquad B=\Z_3[[T]],\qquad R=O[[T]],\qquad
 w=(1+T)^3-1.
\]
The cubic generator fixes all cyclotomic roots.
Complex conjugation $c$ inverts $\sigma$ and commutes with $\gamma$.
Write $\lc(F)$ for the degree-zero logarithmic class group at three.
The logarithmic descent and wild-kernel description in \citet[Definition 1, Proposition 2, and Theorem 3]{JaulentMichel2006} give
\begin{equation}
 \lc(E_n)=M/(\gamma^{3^n}-1)M,\qquad
 U_e^\vee=(\Z_3(e)\otimes M)_{\langle\gamma\rangle}
 =M/(4^e\gamma-1)M\quad(e\geq1).
 \label{eq:Schneider}
\end{equation}
These are compact coinvariants. The sign $+e$ is dual to the Selmer coefficient $\Q_3/\Z_3(-e)$.
The logarithmic coinvariant at weight zero is not an extension of the definition of $\mathcal S_0$.

\begin{proposition}[finite-level genus dimensions]
\label{prop:generators}
The cubic norm annihilates $A_n$, and
\[
 \dim_{\F_3}A_n/\lambda A_n=
 \begin{cases}0&n=0,\\2&n=1,\\7&n\geq2.\end{cases}
\]
In particular $\dim_{\F_3}A_1/3A_1\leq4$.
\end{proposition}
\begin{proof}
Induction with Chevalley's formula gives $\Cl(K_n)_{(3)}=0$:
each cyclotomic step is cyclic cubic with one ramified finite prime, so its invariant three-class group has order at most one; nilpotence on three-torsion kills the whole group.
Norm to $K_n$ therefore makes $1+\sigma+\sigma^2$ vanish on $A_n$.
Since $v_3(269^2-1)=3$, the residue degree of $269$ in $K_n$ is $2\cdot3^{\max(n-2,0)}$.
Hence $E_n/K_n$ has exactly $t_n=3^{\min(n,2)}$ ramified primes, all tame and totally ramified, while every prime at three splits.
Chevalley's formula gives
\begin{equation}
 \#A_n[\lambda]=3^{t_n-1}/j_n,\label{eq:genus}
\end{equation}
where $j_n$ is the unit norm index, an exponent-three group order.

We claim $j_0=j_1=1$ and $j_n=3$ for $n\geq2$.
For a unit $u$ of the CM field $K_n$, $u/\overline u$ is a root of unity by Kronecker's lemma.
Thus $u^2$ is a root of unity times the real unit $u\overline u$; squaring is bijective in the exponent-three norm quotient.
At a prime over $269$, conjugation is the unique order-two residue automorphism.
A real unit reduces to the field of order $269^{3^{\max(n-2,0)}}$, whose multiplicative order is prime to three.
Its residue is a cube, hence it is a norm in the tame cubic local extension; principal units are cubes there.
At the other finite places units are norms in the unramified or split local extensions.
The cyclic Hasse norm theorem \citep[Chapter VIII, Theorem 3.1]{Milne2020} makes every real unit a global norm.
Thus roots of unity supply the only obstruction.
For $n\geq2$ the residue-field cardinality $q$ satisfies $v_3(q-1)=n+1$, so a primitive $3^{n+1}$-st root is not a cube there.
For $n=0,1$ all the roots of unity are cubes, since $v_3(269^2-1)=3$.
This proves the claimed indices; \eqref{eq:genus} gives the dimensions because kernel and cokernel of multiplication by $\lambda$ have equal orders.
The final bound follows from $3O=\lambda^2O$ and the two successive quotients of $A_1/\lambda^2A_1$.
\end{proof}

\begin{proposition}[the ray genus quotient]
\label{prop:raygenus}
The split cyclotomic module satisfies
\[
 M/\delta M\simeq\F_3[T]/(T^7).
\]
It is generated by seven elements over $O$.
\end{proposition}
\begin{proof}
The preceding norm argument makes $M$ an $R$-module.
There are nine primes above $269$ in $K_\infty$; their stabilizer in $\Gamma$ is $\Gamma^9$.
For $n\geq2$, each residue field has cardinality $269^{2\cdot3^{n-2}}$ and its three-primary multiplicative group has order $3^{n+1}$.
Conjugation fixes each of the nine primes and acts by inversion on this group, since
\[
 v_3\bigl(269^{3^{n-2}}+1\bigr)=n+1.
\]

Let $Y_n$ be the three-primary ray class group for the product of the primes above $269$, with the additional condition that the unique prime above three split completely.
The ray exact sequence \citep[Chapter V, Section 1]{Milne2020} uses the group of three-units for this splitting condition.
If $u$ is a three-unit, then $u/c(u)$ has valuation zero at the unique prime above three and at all other finite primes.
It is therefore an ordinary unit; every complex conjugate has absolute value one, so it is a root of unity.
Thus the negative part of the three-adic completion of the three-units is exactly $\mu_{3^{n+1}}$.
The positive part maps to zero in the residue three-groups, on which conjugation is negative.
The class-group term vanishes because $\Cl(K_n)_{(3)}=0$.
Consequently, for $n\geq2$,
\[
 Y_n\simeq(\Z/3^{n+1})^9/\langle(1,\ldots,1)\rangle.
\]
The diagonal generator is a primitive $3^{n+1}$-st root of unity reduced at all nine primes.
Taking norm limits, with these compatible root-of-unity identifications, gives
\begin{equation}
 Y\simeq
 \Z_3(1)[\Gamma/\Gamma^9]\big/
 \Z_3(1)\sum_{g\in\Gamma/\Gamma^9}g.
 \label{eq:raymodule}
\end{equation}
This is the ray module over $K_\infty$ with splitting at three imposed; it is free of $\Z_3$-rank eight.
No further negative relation is introduced by passing from units to three-units.

Let $\mathcal L$ be the extension defining $M$, and put
$P=\Gal(\mathcal L/K_\infty)$ and $\Delta=\Gal(E_\infty/K_\infty)$.
Each inertia subgroup at $269$ in $P$ has order three and maps isomorphically to $\Delta$.
The maximal abelian quotient of $P$ is $Y/3Y$.
Indeed its inertia has exponent three, so it is a quotient of $Y/3Y$.
Conversely the extension defined by $Y/3Y$ contains $E_\infty$, and every one of its nine inertia groups maps isomorphically onto $\Delta$.
It is therefore unramified over $E_\infty$ and split above three, hence lies in $\mathcal L$.

An inertia group gives a section of $P\to\Delta$, so $P=M\rtimes\Delta$.
Its commutator subgroup is $(\sigma-1)M$.
Abelianization therefore gives the $\Gamma$-equivariant exact sequence
\[
 0\longrightarrow M/\delta M\longrightarrow Y/3Y
 \longrightarrow\Delta\longrightarrow0.
\]
Modulo three the cyclotomic twist in \eqref{eq:raymodule} is trivial, and the norm element of the cyclic group of order nine is $T^8$.
Hence $Y/3Y=\F_3[T]/(T^8)$.
The quotient to $\Delta\simeq\F_3$ is nonzero and $\Gamma$ acts trivially on $\Delta$; its kernel is the ideal generated by $T$.
This proves the displayed isomorphism.
Compact Nakayama over $O$ now gives seven generators: lift a basis of $M/\delta M$ and take its closed $O$-span; the quotient equals three times itself and therefore has no nonzero finite quotient.
\end{proof}

For a finite arithmetic check of the ray calculation, the three-primary groups at levels $n=0,1,2$ are respectively
\[
 \Z/9,\qquad (\Z/27)^2\oplus\Z/3,\qquad(\Z/27)^8.
\]
The unique prime above three has trivial three-primary ray class at each level, so imposing its splitting does not change these groups.
The program \emph{RayClassStart}, available through the workbench, computes them from the cyclotomic polynomials and modulus $269$, after complete class, unit, and maximal-order certification.
The all-level assertion in Proposition~\ref{prop:raygenus} follows from the ray exact sequence, independently of this finite check.

\section{Finite arithmetic certificates}
\label{sec:certificates}
This section states the finite arithmetic information and gives the certificates by which it is verified.
\begin{proposition}[certificate package]
\label{prop:certificates}
The following finite arithmetic statements hold:
\begin{enumerate}
\item[(C1)] $\lc(E)\simeq\Z/9\oplus\Z/3$.
\item[(C2)] $\#\lc(E_1)\leq3^6$.
\item[(C3)] For $F=\Q(\alpha)$ and $F_1=\Q(\alpha,\zeta_9+\zeta_9^{-1})$, one has $\lc(F)\simeq\Z/3$ and $\lc(F_1)\simeq\Z/9$.
\item[(C4)] On $\lc(F_1)$, the automorphism fixing $\alpha$ and sending $\zeta_9+\zeta_9^{-1}$ to $\zeta_9^4+\zeta_9^{-4}$ acts as multiplication by four.
\end{enumerate}
The proof is given by the exact witnesses and precision checks below and in Appendix~\ref{app:witnesses}.
\end{proposition}

\subsection{Field identifications and logarithmic computations}
For $E$ use
\begin{equation}
 P(z)=z^6+3z^5+6z^4+545z^3+813z^2-804z+71824.
 \label{eq:polyE}
\end{equation}
The generator $z=\omega-\alpha$ satisfies this polynomial, and
\[
 \alpha=\frac{269-(2z+1)(z^2+z+1)}{3z(z+1)},\qquad\omega=z+\alpha
\]
verify the reverse inclusion exactly.
The certified outputs are ordinary class invariants $[7,7]$, full \texttt{bnfcertify(b)=1}, and
\texttt{bnflog(b,3)=[[9,3],[9,3],[]]}.
The first component is the degree-zero logarithmic group.
For $F$ the polynomial is $z^3-269$, with certified ordinary classes $[7]$, full certification, and logarithmic output \texttt{[[3],[3],[]]}.
For $F_1$ use
\begin{align}
 p(z)={}&z^9-3z^7-270z^6+810z^5+1347z^4\notag\\
 &-48600z^3-72900z^2-96123z-728731.
 \label{eq:polyF1}
\end{align}
The certified ordinary classes are $[21,3]$, with full certification and logarithmic output \texttt{[[9],[3],[3]]}.
Exact embedded roots of $X^3-269$ and $X^3-3X+1$ identify this degree-nine field with $F_1$.
Its signature is $(3,3)$: conjugation-fixed does not mean totally real.
The algorithms and finite stopping criterion are proved by \citet[Sections 2--4, especially Lemma 3.8]{BelabasJaulent2016}.
The class and unit data in these computations have full certification. In the upper-bound calculation below, completeness of a proposed class or unit basis is unnecessary.

\subsection{The degree-eighteen upper certificate}
A full class-number computation for $E_1$ is unnecessary.
Four explicit irreducible pure cubics $X^3-a_i$ have relative discriminant the unit ideal; since $\mu_3\subset E_1$, they define unramified cyclic cubic extensions.
Against four named ideals their cubic residue-symbol matrix is
\begin{equation}
 \begin{pmatrix}0&2&0&0\\0&0&2&0\\2&0&0&0\\0&0&0&2\end{pmatrix}
 \quad\text{over }\F_3,
 \label{eq:Artin}
\end{equation}
of determinant one.
For the explicitly defined ideals $g_1,g_2,g_3,q$ of Appendix~\ref{app:witnesses}, exact principal-ideal identities are
\begin{equation}
 g_1^{63}=(b_1),\qquad g_2^{21}=(b_2),\qquad
 g_3^3=(b_3),\qquad q^3=(b_q).
 \label{eq:idealrelations}
\end{equation}
The tested ideals in \eqref{eq:Artin} are precisely $g_1,g_2,g_3,q$.
Proposition~\ref{prop:generators} bounds $\dim_{\F_3}A_1/3A_1$ by four.
The nondegenerate pairing shows that these four classes generate that quotient.
Nakayama therefore makes their three-primary projections generate $A_1$, and the power relations give $\#A_1\leq3^{2+1+1+1}=3^5$.
For at least one unramified cubic, exact local cubeness tests detect a noncube radical at a prime over three.
Its Frobenius is nontrivial; thus the class of that prime has three-primary order at least three.
It follows that $\#A'_1\leq3^4$.

All three completions at three are $\Q_3(\mu_9)$, with logarithmic degree $d=3\log_3(4)$.
Two explicit $S_3$-units have their first two logarithmic valuation coordinates
\begin{equation}
 \begin{pmatrix}39&51\\21&30\end{pmatrix}\pmod{81}.
 \label{eq:logminor}
\end{equation}
The determinant is $18\pmod{81}$, with exact three-adic valuation two.
The product formula makes the third coordinate minus the sum of the first two.
Thus these units span a sublattice of index nine in the rank-two degree-zero divisor lattice supported above three.
Additional true unit relations can only decrease this index.
The logarithmic kernel supported above three has order at most nine, and its quotient maps into $A'_1$.
Hence $\#\lc(E_1)\leq3^6$.
Surjectivity onto the full $S$-class group is not assumed.

This certificate requires a certified maximal order, exact ideal identities, relative irreducibility and discriminants, residue-symbol and local noncube tests, and absolute precision sufficient for \eqref{eq:logminor}.
It does not assume that provisional class generators or units form complete bases; the pairing and genus bound prove the required generation independently.

\subsection{The labelled cyclotomic action}
In the fully certified degree-nine field, choose ordinary class representatives $\mathfrak a_1,\mathfrak a_2$ coprime to three, of orders $21,3$.
The primes $\mathfrak p_1,\mathfrak p_2$ above three have local degrees six and three and common logarithmic degree $d=3\log_3(4)$; both are fixed by $\gamma$.
Put
\[
 r_i=\frac{\log_3N\mathfrak a_i}{d},\qquad
 v_j(u)=-\frac{\log_3N_{(F_1)_{\mathfrak p_j}/\Q_3}(u)}{d}.
\]
A basis of the degree-zero free divisor module is
\[
 \mathfrak a_1-r_1\mathfrak p_2,\quad
 \mathfrak a_2-r_2\mathfrak p_2,\quad\mathfrak p_1-\mathfrak p_2.
\]
Principal generators for the two class-power relations, two decompositions of the primes in the class basis, and five fundamental units give a complete relation presentation.
Indeed, every other principal relation has the same ordinary divisor as a combination of the four former relations; their quotient is a certified unit.
A finite-index unit subgroup would be insufficient for this completeness argument.
For the representatives used in the explicit certificate, the relation and action matrices are
\begin{equation}
 R_0\equiv\begin{pmatrix}
 21&0&26&8&0&0&0&0&0\\
 0&3&0&0&0&0&0&0&0\\
 6&4&24&21&0&0&18&6&6
 \end{pmatrix},\qquad
 G_0\equiv\begin{pmatrix}1&14&0\\0&1&0\\0&8&1\end{pmatrix}\pmod{27}.
 \label{eq:matrices}
\end{equation}
The primitive row $q_0=(0,1,6)$ has $q_0R_0=0$ and $q_0G_0=4q_0\pmod9$.
The finite Smith invariants of $(R_0\mid27I_3)$ are $[9,1,1]$.
The row therefore identifies the group with $\Z/9$ and computes its action.
The sign in the local correction matters: from $\gamma\mathfrak a_i=(u_i)\prod_k\mathfrak a_k^{b_{ki}}$, the logarithmic principal relation gives
\[
 [\gamma\mathfrak a_i]=\sum_kb_{ki}[\mathfrak a_k]-\sum_jv_j(u_i)[\mathfrak p_j].
\]
For $\beta=\zeta_9+\zeta_9^{-1}$, $\beta^3-3\beta+1=0$ and $\gamma\beta=2-\beta-\beta^2$.
This labels the value-four automorphism, rather than its inverse.
The independent reconstruction identifies this automorphism exactly, determines all primitive quotient rows modulo nine, and repeats the local calculation at two precisions.

\section{Specialization of the cyclotomic module}
This section combines the ray genus quotient and the finite logarithmic data to evaluate every positive weight.
The following elementary lemma isolates the required power-series calculation.

\begin{lemma}[specialization from a first-layer bound]
\label{lem:specialization}
Let $J,H$ be ideals of $B=\Z_3[[T]]$ satisfying
\[
 J\subset H,\qquad 3H\subset J,\qquad
 J+(T)=(9,T),\qquad H+(w)=(9,T-3),
\]
where $w=(1+T)^3-1$.
Suppose
\[
 (J+(3))/(3)=(T^d)\subset\F_3[[T]],\qquad d\geq4,
\]
and
\[
 \#\bigl(B/(J,w)\oplus B/(H,w)\bigr)\leq3^6.
\]
Then $I=J\oplus\delta H$ is an ideal of $R=O[[T]]$, and for every $s\in3\Z_3$,
\begin{equation}
 R/(I,T-s)\simeq
 \begin{cases}
 O/\delta^3,&s\not\equiv3\pmod9,\\
 O/\delta^4,&s\equiv3\pmod9.
 \end{cases}
 \label{eq:specialization}
\end{equation}
Moreover $H=(9,T-3)$ and $R/I$ is finite.
\end{lemma}
\begin{proof}
The inclusions on $J,H$ and the relation $\delta^2=-3$ show that $I$ is stable under multiplication by $\delta$, hence is an ideal of $R$.
Put $C=B/(J,w)$.
The other first-layer summand is $B/(9,T-3)\simeq\Z/9$, so $\#C\leq3^4$.
On the other hand,
\[
 C/3C\simeq\F_3[T]/(T^3),\qquad C/TC\simeq\Z/9.
\]
Thus $C$ has three generators as a three-group and exponent at least nine.
These facts and the order bound force
\[
 C\simeq\Z/9\oplus(\Z/3)^2.
\]
In particular there exist $j\in J$ and $a\in B$ with
\[
 9=j+wa.
\]
Reducing modulo three gives $\overline j=-T^3\overline a$.
The genus hypothesis implies
\[
 a\in(3,T^{d-3})\subset(3,T).
\]
For $s\in3\Z_3$ one has $w(s)\in9\Z_3$ and $a(s)\in3\Z_3$, so
\begin{equation}
 j(s)=9-w(s)a(s)\equiv9\pmod{27}.
 \label{eq:unitnine}
\end{equation}
Every element of $J$ has constant coefficient divisible by nine, and its coefficients of $T^i$ for $1\leq i<d$ are divisible by three.
Its value at any $s\in3\Z_3$ is therefore divisible by nine.
Equation~\eqref{eq:unitnine} proves that the evaluation ideal is exactly
\[
 J(s)=(9)\qquad(s\in3\Z_3).
\]

The inclusion $H\subset(9,T-3)$ shows that $H(s)\subset3\Z_3$ always and $H(s)\subset9\Z_3$ when $s\equiv3\pmod9$.
Lift $T-3$ through $H+(w)$, writing $T-3=h+wb$ with $h\in H$ and $b\in B$.
Then
\[
 h(s)\equiv s-3\pmod9.
\]
If $s\not\equiv3\pmod9$, this value has valuation one and $H(s)=(3)$.
If $s\equiv3\pmod9$, the inclusion $J\subset H$ and $J(s)=(9)$ instead give $H(s)=(9)$.
Since evaluation identifies $R/(I,T-s)$ with $O/(J(s)+\delta H(s))$, the two resulting ideals are $(\delta^3)$ and $(\delta^4)$, proving \eqref{eq:specialization}.

Finally $h=T-3-wb$ has derivative congruent to one modulo three and constant term $-3$.
Hensel's lemma gives a root $c_0\in3\Z_3$, with $c_0\equiv3\pmod9$.
Power-series division writes $h=(T-c_0)u$ for a unit $u\in B$, so $T-c_0\in H$.
Dividing the element $j$ from \eqref{eq:unitnine} by $T-c_0$ shows that $j(c_0)\in H$.
This value is nine times a unit, hence $9\in H$.
Thus $H=(9,T-c_0)=(9,T-3)$.
It follows that $I$ contains $9\delta=\delta^5$.
The successive quotients of the finite filtration of $R/I$ by powers of $\delta$ are quotients of
\[
 R/(I,\delta)=\F_3[T]/(T^d).
\]
They are finite, proving the last assertion.
\end{proof}

\begin{proposition}[the arithmetic presentation]
\label{prop:arithmeticpresentation}
The module $M$ has a cyclic $R$-presentation $M=R/(J\oplus\delta H)$ satisfying Lemma~\ref{lem:specialization} with $d=7$.
Consequently $M$ is finite, its positive component is $\Z/9$ with $\gamma$ acting by four, and
\[
 \lc(E_1)\simeq(\Z/9)^2\oplus(\Z/3)^2
\]
as an abelian group.
\end{proposition}
\begin{proof}
For a quadratic extension, extension and norm on logarithmic divisor classes satisfy $Nj=2$ and $jN=1+c$.
The maps preserve degree zero, by compatibility of local degrees with norms, and commute with $\gamma$.
Since two is invertible over $\Z_3$, they identify the downstairs logarithmic group with the upstairs positive subgroup:
\[
 \lc(F)=\lc(E)^{c=1},\qquad\lc(F_1)=\lc(E_1)^{c=1}.
\]
By (C1), $M/TM$ has additive type $\Z/9\oplus\Z/3$.
The only $O$-module of length three with an element of order nine is $O/\delta^3$.
Its positive component has order three by (C3), so a cyclic generator has negative sign.
Lift such a generator and apply $(1-c)/2$ to obtain an odd element $x\in M$.
Compact Nakayama makes $x$ a cyclic $R$-generator.

The annihilator of $x$ is invariant under coefficient conjugation.
Since $R=B\oplus\delta B$ and two is invertible, it is of the form $J\oplus\delta H$ for ideals $J,H\subset B$.
Stability under $\delta$ gives $J\subset H$ and $3H\subset J$.
Because $x$ is odd,
\[
 M^-=B/J,\qquad M^+=B/H.
\]
The base negative component is $\Z/9$, giving $J+(T)=(9,T)$.
The positive component of $M/wM$ is the logarithmic group of $F_1$.
By (C3)--(C4), it is cyclic of order nine with $T$ acting by three, so
\[
 H+(w)=(9,T-3).
\]
Moreover $M/\delta M=B/(J,3)$, and Proposition~\ref{prop:raygenus} gives $J\bmod3=(T^7)$.
The first-layer order bound is exactly (C2).
All hypotheses of Lemma~\ref{lem:specialization} hold.
Its proof also identifies the negative first-layer component with $\Z/9\oplus(\Z/3)^2$; adjoining the positive component $\Z/9$ gives the stated logarithmic group.
\end{proof}

\begin{proof}[Proof of Theorem~\ref{thm:main}]
In \eqref{eq:Schneider}, the relation $4^e\gamma=1$ is $T=4^{-e}-1\in3\Z_3$.
Modulo nine, this value is $-3e$.
Lemma~\ref{lem:specialization} and Proposition~\ref{prop:arithmeticpresentation} therefore give
\[
 U_e^\vee\simeq O/\lambda^{3+\ind_{e\equiv2\ (3)}}.
\]
Proposition~\ref{prop:structure} adds the local length $2(1+v_3(e))$ and gives cyclicity of the entire $K$-group, proving the formula for $\kappa_e$.
Proposition~\ref{prop:descent} adds one to obtain $m_e$.
Proposition~\ref{prop:classes} and Lemma~\ref{lem:residual} give the vanishing assertions.
There are $\lfloor(N+1)/3\rfloor$ positive integers at most $N$ congruent to two modulo three, proving the sum.
\end{proof}

\Needspace{18\baselineskip}
\section{Explicit groups and scope}
This section translates the module formula into abelian groups and records the scope of the result.
The underlying abelian groups satisfy
\[
 O/\lambda^{2h}\simeq(\Z/3^h)^2,\qquad
 O/\lambda^{2h+1}\simeq\Z/3^{h+1}\oplus\Z/3^h.
\]
For example, $e=1$ has strict length $a_1=3$ and local correction $2$, so $\kappa_1=5$.
The odd-length formula gives $K_2(\Oa_E)_{(3)}\simeq\Z/27\oplus\Z/9$, of order $3^5$.
The first weights give the following abelian groups:
\begin{center}\begin{tabular}{@{}rrrrl@{}}\toprule
$e$&$\kappa_e$&$a_e$&$m_e$&$K_{2e}(\Oa_E)_{(3)}$\\\midrule
1&5&3&4&$\Z/27\oplus\Z/9$\\
2&6&4&5&$\Z/27\oplus\Z/27$\\
3&7&3&4&$\Z/81\oplus\Z/27$\\
4&5&3&4&$\Z/27\oplus\Z/9$\\
5&6&4&5&$\Z/27\oplus\Z/27$\\
6&7&3&4&$\Z/81\oplus\Z/27$\\
9&9&3&4&$\Z/243\oplus\Z/81$\\\bottomrule
\end{tabular}\end{center}
The strict character lengths are periodic modulo three, whereas the full even $K$-groups also contain the unbounded local term $2v_3(e)$.
Ordinary class-group vanishing fixes the residual rank, but cannot alone bound the exponent.
The first-layer bound and labelled cyclotomic action supply the necessary additional information.
The $K$-theoretic conclusions concern one field and one coefficient prime. They do not estimate all coefficient primes or varying host fields and do not prove a Riemann hypothesis for the associated cubic Hecke function.

\subsection{The separate seven-primary result}

\begin{proof}[Proof of Theorem~\ref{thm:seven-stability}]
Appendix~\ref{seven:section} gives a full finite certification of
$\Cl(\Q(\alpha))_{(7)}=\Cl(F_1^{(7)})_{(7)}=\Z/7\Z$:
explicit independent units, a nonsingular matrix of seventh-power residue
characters, and bounded $L$-value and regulator computations isolate each
class-number-to-unit-index ratio as seven and show that its index is prime
to seven. Every exact algebraic witness and every truncation bound is
printed there.
The prime seven is inert in $\Q(\alpha)$, since $X^3-3$ is irreducible over
$\F_7$, and is the unique ramified prime in its cyclotomic tower.
Norm coinvariants therefore identify with the lower class group.
Equality in the first two layers and Nakayama's lemma prove constancy in
all layers, the one-prime case of \citet[Theorem~1(1)]{sevenFukuda1994}.
Projection of the conjugation-fixed class group onto either faithful
cubic-character component gives the displayed inverse-limit module.
The classes of seven-primes are $\Delta$-trivial, and
$\psi(\delta)-1$ is a seven-adic unit for a generator $\delta$ of $\Delta$.
Global and local inflation--restriction consequently preserve the strict
conditions. Equivariance specializes the module
$\Z_7[[T_7]]/(7,T_7)$ at $T_7=\kappa(\gamma_7)-1$, giving the stated quotient.
The appendix proves these finite control and character-projection
steps in full.
\end{proof}

\section{Real nonvanishing of the associated cubic Hecke function}
\label{sec:cubic-real-nonvanishing}
This section proves that the completed cubic function has no real zero, proves that the uncompleted function has no positive real zero, and supplies the Mellin formula used in the class-number certificate.
The argument adapts the finite partial-sum positivity method of \citet[Theorem~4, pp.~646--647]{LouboutinMunsch2013}, which they formulate for theta series of real even Dirichlet characters.
Their kernel has square exponents and coefficients bounded by one; here the exponents are linear and the coefficients have a divisor bound. The proof below establishes this variation directly.
The conductor-specific finite arithmetic is checked by two coefficient algorithms and by a separate computation recorded in the workbench.

Put $F=\mathbb Q(\alpha)$, $\alpha^3=269$, and retain $K=\mathbb Q(\omega)$ and $E=FK$.
Let $\rho$ be the two-dimensional sum-zero representation of $\operatorname{Gal}(E/\mathbb Q)\simeq S_3$.
Write
\[
 L_F^{\mathrm{rel}}(s)=\frac{\zeta_F(s)}{\zeta(s)}=L(s,\rho)
   =\sum_{n\ge1}A(n)n^{-s},\qquad
 Q=3\cdot269^2=217083,\qquad a=\frac{2\pi}{\sqrt Q}.
\]
The superscript distinguishes this complex function from coefficient fields used elsewhere.

\begin{proof}[Proof of Theorem~\ref{thm:cubic-real-positive}]
We first identify the analytic function and its coefficients.
At 269 the pure cubic polynomial is Eisenstein, giving tame ramification of degree three and different exponent two.
At three, $269$ is a cube in $\mathbb Q_3$: its negative lies in $1+9\mathbb Z_3$, and the $3$-adic logarithm identifies the cube map from $1+3\mathbb Z_3$ onto $1+9\mathbb Z_3$.
Thus $F\otimes\mathbb Q_3=\mathbb Q_3\times\mathbb Q_3(\omega)$, with discriminant exponent one.
The polynomial discriminant excludes ramification elsewhere, and the signature is $(1,1)$, so $d_F=-Q$.
The cubic character $\psi$ of $E/K$ has conductor $(269)$: it is tamely ramified there, split at three, and unramified elsewhere.
Induction gives $L_F^{\mathrm{rel}}=L_K(s,\psi)$, which is entire by Hecke's theorem for nontrivial finite-order characters. We use the complete adelic proof in \citet[Theorems~5.16 and~5.22, pp.~34--39]{Poonen2015}; the norm twists of our nontrivial finite cubic character never equal either exceptional norm character. The Artin--Hecke identification is described in \citet[Chapter~VIII, Section~10]{Milne2020}.

At a good rational prime let
\[
 t(p)=\#\{x\bmod p:x^3=269\}-1.
\]
The complete local generating functions are
\[
\begin{array}{c|c}
p&\displaystyle\sum_{k\ge0}A(p^k)T^k\\\hline
3&(1-T)^{-1}\\
269&1\\
p\equiv2\pmod3,\ p\ne269&(1-T^2)^{-1}\\
p\equiv1\pmod3&(1-t(p)T+T^2)^{-1}.
\end{array}
\]
In the last row $t(p)=2$ or $-1$ according as 269 is a cube or a noncube modulo $p$.
The local coefficients are respectively $k+1$ or the periodic sequence $1,-1,0$.
Multiplicativity therefore gives $|A(n)|\le d(n)\le n$.

For an independent coefficient construction, use that $\mathbb Z[\omega]$ has class number one and 269 is inert.
The ray-class exact sequence gives
\[
 \operatorname{Cl}_{(269)}(K)
   \simeq\mathbb F_{269^2}^{\times}/\mu_6,
\]
a cyclic group of order 12060; see \citet[Chapter~V, Section~1]{Milne2020}.
There are exactly two nontrivial cubic characters of this group, which are conjugate.
The character $\psi$ is one of them by the conductor calculation.
For a nonzero residue $v$, identify $v^{24120}$ with a complex cube root of unity and call it $\xi(v)$; set $\xi(0)=0$.
It kills the six global units and is inverted by conjugation.
Consequently
\[
 A(n)=\frac1{12}\sum_{r^2-rs+s^2=n}
   \bigl(\xi(r+s\omega)+\overline{\xi(r+s\omega)}\bigr).
\]
The summand is $2,-1$, or $0$, according as the residue is a nonzero cube, a noncube, or zero.
The factor 12 accounts for six associates and the two conjugate characters, whose norm coefficients coincide.
At three the identity $(1-\omega)^2=-3\omega$ gives $\xi(1-\omega)=1$, agreeing with the Euler factor above.

The \emph{CubicThetaCertificate} computes the Euler coefficients and the lattice coefficients independently, with equality through 8192.
For $S_k=\sum_{n\le k}A(n)$ it verifies the following exact integer assertions:
\[
 \min_{1\le k\le2048}S_k=1,\qquad
 \{k\le2048:S_k=1\}=\{1,2\},\qquad S_{2048}=47.
\]
The first thirteen coefficients, for example, are
\[
 1,0,1,1,0,0,-1,0,1,0,0,1,-1.
\]
Put $N=2048$ and $x=e^{-at}$ with $t\ge1$.
Since $\sqrt{217083}<466$ and $\pi>25/8$, we have $a>1/75$, and $e^{1/75}>1+1/75$ gives $0<x<75/76$.
Abel summation yields
\[
 \sum_{n=1}^N A(n)x^n
   =(1-x)\sum_{k=1}^{N-1}S_kx^k+S_Nx^N\ge x.
\]
The absolute tail divided by $x$ is bounded by
\[
 \sum_{n>N}n x^{n-1}
  \le76(N+76)(75/76)^N<10^{-6},
\]
where the last step is the exact integer comparison
\[
 10^6\cdot76\cdot2124\cdot75^{2048}<76^{2048}.
\]
This proves the stated theta inequality.

Finally, the signatures of $F$ and $\mathbb Q$ are $(1,1)$ and $(1,0)$.
Dividing their completed Dedekind zeta functional equations cancels the real gamma factor and gives
\[
 \Lambda_F^{\mathrm{rel}}(s)=\Lambda_F^{\mathrm{rel}}(1-s)
\]
with sign $+1$.
This uses the Dedekind functional equations and Hecke entireness, rather than an assumption about weight-one modularity.
Let $f(t)=\sum A(n)e^{-ant}$.
Mellin inversion and the functional equation give $f(t)=t^{-1}f(1/t)$ and hence
\begin{equation}\label{eq:cubic-split-mellin}
 \frac{\Lambda_F^{\mathrm{rel}}(s)}2
  =\int_1^\infty f(t)\bigl(t^{s-1}+t^{-s}\bigr)\,dt.
\end{equation}
To justify inversion, shift the Mellin contour in a fixed vertical strip; the completed Hecke function is entire, its gamma factor decays exponentially, and the remaining factor has polynomial growth.
The integral in \eqref{eq:cubic-split-mellin} converges locally uniformly for every complex $s$.
For real $s$ its integrand is strictly positive.
This proves the theorem.
\end{proof}

At $s=1$, termwise integration in \eqref{eq:cubic-split-mellin} gives the formula needed below:
\begin{equation}\label{eq:cubic-L-one}
 L_F^{\mathrm{rel}}(1)
 =\sum_{n\ge1}A(n)\left(\frac{e^{-an}}n+aE_1(an)\right),
 \qquad E_1(u)=\int_u^\infty\frac{e^{-v}}v\,dv.
\end{equation}
Because $E_1(u)\le e^{-u}/u$ and $|A(n)|\le n$, the absolute tail after $M$ terms is at most
\[
 \frac{2e^{-a(M+1)}}{1-e^{-a}}.
\]
This is a rigorous uniform error bound for the interval class-number computation; it requires no zero-free hypothesis.

\appendix
\section{Exact witnesses for the first-layer bound}
\label{app:witnesses}
This appendix gives every algebraic input for the bound $\#\lc(E_1)\leq3^6$ in a compact tower basis.
Put $\zeta=\zeta_9$, so $\zeta^6+\zeta^3+1=0$, and retain $\alpha^3=269$.
The eighteen elements $\alpha^i\zeta^j$, $0\leq i\leq2$, $0\leq j\leq5$, form a rational basis of $E_1$.
For each element $v$ in the tables, define its coefficients by
\begin{equation}
 v=\frac13\sum_{i=0}^2\sum_{j=0}^5 n_{ij}(v)\alpha^i\zeta^j.
 \label{eq:towercoordinates}
\end{equation}
The four radicals are $a_1,\ldots,a_4$; the principal generators in \eqref{eq:idealrelations} are $b_1,b_2,b_3,b_q$; and the two logarithmic-unit witnesses are $u_1,u_2$.

For rational integers $p,t$, write
\[
 \mathfrak p(p,t)=(p,t-\zeta^2-\zeta^5-\alpha\zeta^3)\Oa_{E_1}.
\]
Then the ideals used in the proof are
\[
 g_1=\mathfrak p(1531,492),\qquad
 g_2=\mathfrak p(1531,52)g_1^2,\qquad
 g_3=\mathfrak p(10779577093,2237331723),
\]
with norms $1531,1531^3,10779577093$, respectively.
The remaining ideal is $q=(7,V/3)\Oa_{E_1}$, of norm $343$, where
\begin{align*}
V={}&-6-14\zeta-3\zeta^3-7\zeta^4+7\zeta^5\\
 &+\alpha(-5\zeta-14\zeta^2-4\zeta^4-14\zeta^5)
 +\alpha^2(-2\zeta^2-3\zeta^5).
\end{align*}

\Needspace{12\baselineskip}
\subsection{Radicals}
The following integer coefficients specify the radicals through \eqref{eq:towercoordinates}.
{\small\setlength{\tabcolsep}{7pt}
\begin{longtable}{@{}lrrrr@{}}
\toprule
$v$&$j$&$n_{0j}(v)$&$n_{1j}(v)$&$n_{2j}(v)$\\\midrule
\endfirsthead
\toprule
$v$&$j$&$n_{0j}(v)$&$n_{1j}(v)$&$n_{2j}(v)$\\\midrule
\endhead
$a_1$&0&740549184998&-111559636254&-489466652\\*
$a_1$&1&526461578010&47651376012&-20015464503\\*
$a_1$&2&-557710995894&128108785755&-6461831178\\*
$a_1$&3&20396613524&-114719297449&17281827501\\*
$a_1$&4&834066208629&-81554748855&-7381727682\\*
$a_1$&5&269271546960&86395631373&-19845474474\\
\addlinespace[3pt]
$a_2$&0&17889212813&1549834439&-545756019\\*
$a_2$&1&-19864230075&3173853810&-306758616\\*
$a_2$&2&-15763020834&950364843&655799310\\*
$a_2$&3&19142329988&-3329148144&-326501774\\*
$a_2$&4&11999829114&2955723279&-510496851\\*
$a_2$&5&-22253219370&3228000099&-25475754\\
\addlinespace[3pt]
$a_3$&0&8594118&2058113&-525067\\*
$a_3$&1&-1137294&-2872287&472221\\*
$a_3$&2&-6432990&3329397&-361389\\*
$a_3$&3&21889964&-1332190&-318918\\*
$a_3$&4&-19670253&176235&444729\\*
$a_3$&5&15061875&997560&-515994\\
\addlinespace[3pt]
$a_4$&0&-3061198634644&92343032945&59336272011\\*
$a_4$&1&212898960429&472122029658&-78504530988\\*
$a_4$&2&2196733098180&-302204888193&-6477488079\\*
$a_4$&3&-2778672713688&-4700493766&66964967012\\*
$a_4$&4&-476947997430&207271989861&-21302187252\\*
$a_4$&5&-740558314092&219766407735&-16807475550\\
\addlinespace[3pt]
\bottomrule
\end{longtable}}
\Needspace{12\baselineskip}
\subsection{Principal generators}
The following integer coefficients specify the principal generators through \eqref{eq:towercoordinates}.
{\small\setlength{\tabcolsep}{7pt}
\begin{longtable}{@{}lrrrr@{}}
\toprule
$v$&$j$&$n_{0j}(v)$&$n_{1j}(v)$&$n_{2j}(v)$\\\midrule
\endfirsthead
\toprule
$v$&$j$&$n_{0j}(v)$&$n_{1j}(v)$&$n_{2j}(v)$\\\midrule
\endhead
$b_1$&0&-934103308941123&30779733833058&18279785822499\\*
$b_1$&1&-238943503101082&-48556582298655&26610955045393\\*
$b_1$&2&180507700466093&-131446987179252&19698995146918\\*
$b_1$&3&-351741143032598&-133645627606489&16527571036303\\*
$b_1$&4&566554293146475&-147392907765338&12959157774388\\*
$b_1$&5&1031068311570735&-79098648409229&100889473147\\
\addlinespace[3pt]
$b_2$&0&349881926219204&53877109777038&8241098024896\\*
$b_2$&1&-123303059157588&-18218450949555&-2954745153321\\*
$b_2$&2&-122058690498127&-18289991796469&-2773900234521\\*
$b_2$&3&694489002750061&106586791694733&16559353821251\\*
$b_2$&4&-648097745079781&-100363008485051&-15685611709540\\*
$b_2$&5&526884043868563&82625182552039&12704262121308\\
\addlinespace[3pt]
$b_3$&0&-594138651063&82865774067&1420985001\\*
$b_3$&1&-373650200471&-43157292937&15652197322\\*
$b_3$&2&464371298539&-97854144574&4014966085\\*
$b_3$&3&-59213991169&92038679224&-12836829232\\*
$b_3$&4&-652244114375&57882568157&6685544314\\*
$b_3$&5&-167308011966&-71936274198&15158694417\\
\addlinespace[3pt]
$b_q$&0&5854&845&142\\*
$b_q$&1&4516&660&92\\*
$b_q$&2&-4433&-701&-93\\*
$b_q$&3&456&-60&3\\*
$b_q$&4&6382&1036&156\\*
$b_q$&5&2204&281&63\\
\addlinespace[3pt]
\bottomrule
\end{longtable}}
\Needspace{12\baselineskip}
\subsection{Logarithmic units}
The following integer coefficients specify the logarithmic units through \eqref{eq:towercoordinates}.
{\small\setlength{\tabcolsep}{7pt}
\begin{longtable}{@{}lrrrr@{}}
\toprule
$v$&$j$&$n_{0j}(v)$&$n_{1j}(v)$&$n_{2j}(v)$\\\midrule
\endfirsthead
\toprule
$v$&$j$&$n_{0j}(v)$&$n_{1j}(v)$&$n_{2j}(v)$\\\midrule
\endhead
$u_1$&0&144359&-33486&1723\\*
$u_1$&1&251883&-14043&-3870\\*
$u_1$&2&5356&20922&-3370\\*
$u_1$&3&-71779&-22363&5187\\*
$u_1$&4&161242&-39031&2176\\*
$u_1$&5&140406&-823&-3241\\
\addlinespace[3pt]
$u_2$&0&196813&30434&4705\\*
$u_2$&1&-36473&-5644&-875\\*
$u_2$&2&-208687&-32405&-5026\\*
$u_2$&3&160703&24772&3845\\*
$u_2$&4&160057&24809&3841\\*
$u_2$&5&-104336&-16282&-2510\\
\addlinespace[3pt]
\bottomrule
\end{longtable}}
\subsection{Exact verification procedure}
The verification begins with the defining polynomial
\begin{align*}
P_1(y)={}&y^{18}+1617y^{15}+1082193y^{12}+385574654y^9\\
 &+77580889323y^6+8385564626034y^3+383147566015591.
\end{align*}
It is the compositum polynomial for $X^3-269$ and $X^6+X^3+1$.
An integral basis is computed from $P_1$ and its maximality is certified; in PARI/GP the assertion is \texttt{nfcertify(nfinit(P1))==[]}.
Exact embedded roots of the two constituent polynomials give $\alpha,\zeta$ and the basis above.
Any choices of these roots give an isomorphic certificate because $E_1/\Q$ is Galois and all the displayed elements are transported together.

For each $i$, factor $X^3-a_i$ over $E_1$ and compute its relative discriminant.
The factorization has one degree-three factor and the discriminant is the unit ideal in all four cases.
For a prime $\mathfrak l$ dividing a test ideal, compute $a_i^{(N\mathfrak l-1)/3}$ in its residue field and express the answer as $1,\zeta^3,$ or $\zeta^6$.
Add these exponents with their multiplicities in the test ideal.
This gives exactly \eqref{eq:Artin}.
Exact Hermite forms verify every equality in \eqref{eq:idealrelations}.
The radical $a_1$ is not a cube in any of the three completions above three, as an exact local-power test verifies.
The principal ideals of $u_1,u_2$ have support only at these primes.
These are finite algebraic checks; none requires a complete class-group computation for $E_1$.

Finally factor $P_1$ over $\Q_3$ into its three monic degree-six factors $f_1,f_2,f_3$.
For each unit $u_j$, evaluate its three local logarithmic valuations as
\[
 -\frac{\log_3\operatorname{Res}(f_i,u_j)}{3\log_3(4)}.
\]
Here $u_j$ is expressed as a polynomial in $y$ by the exact field identification.
At both absolute input precisions $40$ and $70$, the matrix modulo $81$, up to permuting its rows, is
\[
 \begin{pmatrix}21&0\\39&51\\21&30\end{pmatrix}\pmod{81}.
\]
The sum of each column vanishes to the computed absolute precision, and every two-row minor has valuation two.
The calculation checks an absolute precision greater than the input precision minus fifteen after division by $3\log_3(4)$, so the valuation two is determined exactly.
This proves the logarithmic index bound used in the main text.

\section{Complete specification of the oriented action calculation}
\label{app:action}
This appendix specifies the complete logarithmic divisor computation that determines the scalar in (C4).
The two arithmetic programs, \emph{CyclotomicAction} and \emph{ExactLogarithmicBound}, are included with their evidence in the accompanying workbench, whose address is given in the method note.
The following PARI/GP code is the complete oriented-action calculation; its only input is the polynomial in \eqref{eq:polyF1}.
The function \texttt{assert} terminates the calculation if a required equality fails.
Complete certification is performed before the class relations and fundamental units are used.
No matrices from \eqref{eq:matrices} are inputs to the program.

\begin{lstlisting}
default(parisize,500000000);
default(parisizemax,2000000000);
default(realprecision,100);
assert(test,label)={if(!test,error(Str("FAILED: ",label)));print("PASS ",label)};
classdecomp(J)={my(z=bnfisprincipal(b,J,1),u);u=nfbasistoalg(b,z[2])/prod(k=1,ng,class_shift[k]^z[1][k]);assert(idealmul(b,idealfactorback(b,class_ideals,z[1]),u)==idealhnf(b,J),"exact ideal decomposition");[z[1],u]};
reconstruct(N)={
 my(fac=factorpadic(p,3,N)[,1],fp,degree,logval,r,cols,rel,act,K,small,q,qn,k,good=0);
 fp=vector(2,j,select(t->poldegree(t)==pr[j].e*pr[j].f,fac)[1]);
 degree=3*log(4+O(3^N));
 logval=(u->(vector(2,j,-log(polresultant(fp[j],lift(nfbasistoalg(b,u))))/degree)~));
 r=vector(ng,j,log(idealnorm(b,class_ideals[j])+O(3^N))/degree);
 cols=concat(concat(vector(ng,j,relpow[j][2]),vector(2,j,relpr[j][2])),b.fu);
 rel=matrix(4,#cols,i,j,if(i>ng,logval(cols[j])[i-ng],if(j<=ng,if(i==j,b.cyc[j],0),if(j<=ng+2,-relpr[j-ng][1][i],0))));
 assert(padicprec(rel,3)>N-20,"relation precision after logarithms");
 assert(vecmin(vector(#cols,j,valuation(sum(i=1,ng,r[i]*rel[i,j])+rel[3,j]+rel[4,j],3)))>N-20,"product formula precision");
 act=matid(4);
 for(j=1,ng,for(i=1,ng,act[i,j]=relgamma[j][1][i]);for(i=1,2,act[ng+i,j]=-logval(relgamma[j][2])[i]));
 K=matrix(4,3,i,j,if(i<=3,if(i==j,1,0),if(j<=ng,-r[j],-1)));
 small=(act*K)[1..3,]; rel=rel[1..3,];
 assert(padicprec(small,3)>N-20,"action precision after logarithms");
 print("PRECISION ",N," RELATIONS_MOD27 ",lift(Mod(rel,27))," ACTION_MOD27 ",lift(Mod(small,27)));
 assert(matsnf(mathnf(concat(lift(Mod(rel,27)),27*matid(3))))==[9,1,1],"finite Smith invariants");
 for(qn=1,728,q=[qn%9,(qn\9)%9,(qn\81)%9];if(gcd(gcd(q[1],q[2]),gcd(q[3],3))==1 && lift(Mod(q*rel,9))==vector(#cols),good++;assert(lift(Mod(q*small-4*q,9))==[0,0,0],"oriented quotient action four")));
 assert(good==6,"all six primitive characters checked");
 [lift(Mod(rel,27)),lift(Mod(small,27))]
};
main()={
p=y^9-3*y^7-270*y^6+810*y^5+1347*y^4-48600*y^3-72900*y^2-96123*y-728731;
assert(polisirreducible(p),"degree-nine field irreducible");
b=bnfinit(p,1);
print("SIGNATURE ",b.nf.sign," DISCRIMINANT ",b.disc," CLASS_INVARIANTS ",b.cyc);
assert(bnfcertify(b)==1,"full class-group and unit certification");

assert(nfcertify(b)==[],"maximal order");
print("LOG_CLASS ",bnflog(b,3));
assert(bnflog(b,3)[1]==[9],"degree-zero logarithmic group cyclic nine");
a=nfroots(b,x^3-269)[1]; beta=nfroots(b,x^3-3*x+1)[1];
assert(a^3==269 && beta^3-3*beta+1==0,"exact embedded cubic fields");
auts=nfgaloisconj(b);
chosen=select(t->nfgaloisapply(b,t,a)==a && nfgaloisapply(b,t,beta)==2-beta-beta^2,auts);
assert(#chosen==1,"unique cyclotomic automorphism of value four");
gamma_aut=chosen[1];
pr=idealprimedec(b,3);
if(pr[1].e*pr[1].f<pr[2].e*pr[2].f,pr=Vecrev(pr));
assert(#pr==2 && pr[1].e*pr[1].f==6 && pr[2].e*pr[2].f==3,"local degrees six and three");
assert(vector(2,j,bnflogef(b,pr[j]))==[[2,3],[1,3]],"logarithmic local indices");
for(j=1,2,assert(idealhnf(b,nfgaloisapply(b,gamma_aut,pr[j]))==idealhnf(b,pr[j]),Str("prime fixed ",j)));
ng=#b.gen;
assert(ng==2 && b.cyc==[21,3],"ordinary class presentation");
class_shift=vector(ng,j,nfbasistoalg(b,idealcoprime(b,b.gen[j],3)));
class_ideals=vector(ng,j,idealmul(b,b.gen[j],class_shift[j]));
assert(matrix(2,ng,j,i,idealval(b,class_ideals[i],pr[j]))==matrix(2,ng),"representatives coprime to three");
relpow=vector(ng,j,classdecomp(idealpow(b,class_ideals[j],b.cyc[j])));
relpr=vector(2,j,classdecomp(pr[j]));
relgamma=vector(ng,j,classdecomp(nfgaloisapply(b,gamma_aut,class_ideals[j])));
for(j=1,ng,assert(relpow[j][1]==Col(vector(ng)),Str("principal class power ",j)));
R60=reconstruct(60);R90=reconstruct(90);assert(R60==R90,"local precision agreement");
print("ACTION_CERTIFIED");
};
main();
quit
\end{lstlisting}

\section{Ordinary class groups by ideal relations and interval bounds}
This appendix determines the ordinary class groups using integer ideal arithmetic and certified analytic bounds.
Let $F=\mathbb Q(\alpha)$ with $\alpha^3=269$, and let $E=F(\omega)$.
It also verifies a fundamental unit without depending on the logarithmic class-group computations.
\begin{proposition}
The ordinary ideal class groups are
\[
 \operatorname{Cl}(F)\simeq C_7,\qquad
 \operatorname{Cl}(E)\simeq C_7\oplus C_7.
\]
The element $\epsilon=108\alpha^2-739\alpha+270$ generates the free part of $\mathcal O_F^\times$.
\end{proposition}
\begin{proof}
An exact integral-basis computation gives
\[
 \mathcal O_F=\mathbb Z\left[1,\alpha,
 \beta=\frac{1+2\alpha+\alpha^2}{3}\right],\qquad
 d_F=-217083.
\]
Minkowski's ideal-class bound \citep[Theorem~4.3]{MilneANT2020} is less than $132$.
Indeed, the inequalities $\pi>157/50$ and $\sqrt{217083}<466$ give
\[
 \frac4\pi\frac6{27}\sqrt{217083}<132.
\]
The ideal classes of prime ideals of norm at most $131$ therefore generate the class group.
We use the forty-eight prime ideals above rational primes at most $131$ obtained by omitting inert primes, whose prime ideals are principal; including conjugate prime ideals of larger norm only enlarges this generating set.
The rational primes give twenty-two principal relations.
A further twenty-seven relations are furnished by $a+b\alpha+c\beta$, where $(a,b,c)$ runs through
\[
\begin{gathered}
(-2,-2,1),(-2,-2,2),(-2,-1,1),(-2,1,0),(-2,2,0),\\
(-2,2,2),(-1,1,2),(0,-2,1),(0,1,1),(1,-1,2),\\
(1,1,1),(1,1,2),(2,0,1),(2,2,1),(-4,-4,3),\\
(-4,-3,2),(-4,-2,1),(-4,1,0),(-4,1,1),(-4,1,3),\\
(-4,4,1),(-3,-3,2),(-3,-2,1),(-3,-2,2),(-1,-4,1),\\
(3,-2,1),(-5,1,4).
\end{gathered}
\]
Each relation is verified as follows. For $u=3a+c$, $v=3b+2c$, and $w=c$, the norm is the integer
\[
 N_{F/\mathbb Q}(a+b\alpha+c\beta)
 =\frac{u^3+269v^3+269^2w^3-807uvw}{27}.
\]
Factoring this integer gives the possible rational prime divisors. Exact lattice membership in each successive power of each prime ideal gives its valuation. The product of the resulting ideal norms equals the absolute element norm, proving that the factorization is complete. The Hermite normal form of the $48$ by $49$ relation matrix has determinant $7$. Hence $h_F\leq7$.

For a lower bound, use $L(s)=L_F^{\mathrm{rel}}(s)$ and $a=2\pi/\sqrt{217083}$ from Theorem~\ref{thm:cubic-real-positive}.
The coefficients are determined by the local Euler factors in Section~\ref{sec:cubic-real-nonvanishing}, and satisfy $|A(n)|\leq d(n)\leq n$.
The split Mellin formula is \eqref{eq:cubic-L-one}.
Since $aE_1(an)\leq e^{-an}/n$, truncation after $8192$ terms has absolute error at most
$2e^{-8193a}/(1-e^{-a})<1.551\cdot10^{-46}$.
Directed interval arithmetic at $220$ bits, with all coefficients independently computed by root counting, proves
\[
\begin{split}
L(1)&\in[1.6976757232625543059851766859761711810674471956,\notag\\[-2pt]
 &\hspace{16mm}1.6976757232625543059851766859761711810674471964],\\
R_\epsilon=-\log|\epsilon|&\in[17.98414259051485567560205694050578826067542238514,\notag\\[-2pt]
 &\hspace{16mm}17.98414259051485567560205694050578826067542238515].
\end{split}
\]
Here the real embedding is understood, and the norm of $\epsilon$ is exactly $1$.
In particular
\[
 6<\frac{\sqrt{217083}\,L(1)}{2\pi R_\epsilon}<8.
\]
If $m$ is the index of $\langle-1,\epsilon\rangle$ in $\mathcal O_F^\times$, the analytic class-number formula \citep[Chapter~V, Theorem~2.4]{Milne2020} identifies this ratio with $h_F/m$. Since $h_F\leq7$ and $m$ is a positive integer, the strict lower bound $6$ forces $m=1$ and $h_F=7$.

We now deduce the class group of $E$ without a second degree-six class-group algorithm.
Let $\sigma(\alpha)=\omega\alpha$ and let $c$ be complex conjugation.
Every cyclotomic subfield of $E$ is contained in its maximal abelian subfield $\mathbb Q(\omega)$, so the roots of unity of $E$ form $\mu_6$.
The free unit lattice
$M_u=\mathcal O_E^\times/\mu_6$ has rank two over $\mathbb Z$.
The operator $1+\sigma+\sigma^2$ is zero on $M_u$, since its norm takes values in the finite unit group of $\mathbb Q(\omega)$.
Thus $M_u$ is a torsion-free rank-one module over the principal ideal domain $\mathbb Z[\omega]$, and is free.
In a suitable generator, $c$ acts as $z\mapsto\overline z$ or $z\mapsto-\overline z$: its initial form is $z\mapsto u\overline z$, where $u\in\mu_6$, and changing the generator changes $u$ by a square in $\mu_6$.
Consequently
\[
 q=[M_u:M_u^c+\sigma M_u^c]\in\{1,3\}.
\]
Indeed the two fixed lattices are respectively $\mathbb Z$ and
$\mathbb Z\sqrt{-3}$, whose $\mathbb Z[\omega]$-spans have these indices.

The image of $\mathcal O_F^\times$ in $M_u$ equals $M_u^c$.
To prove the nontrivial inclusion, choose a unit $u$ representing a fixed class.
Then $c(u)/u\in\mu_6$, and multiplication of $u$ by a root of unity reduces this ratio to $1$ or $-1$.
The second case would give $u=\sqrt{-3}\,v$ for $v\in F$, and
$N_{E/F}(u)=3v^2$ would be a unit of $F$.
But $F$ has a prime above $3$ with ramification index one, since $269$ is a cube in $\mathbb Q_3$.
Its valuation in $3v^2$ would be odd, a contradiction.
Thus $u$ can be chosen in $F$.

At the three complex places of $E$, the logarithm vectors of $\epsilon$ and $\sigma\epsilon$, with the usual factor two at complex places, are permutations of
$(-2R_\epsilon,R_\epsilon,R_\epsilon)$ and
$(R_\epsilon,-2R_\epsilon,R_\epsilon)$.
Their regulator is $3R_F^2$, so
$R_E=3R_F^2/q$.
Artin formalism gives
\[
 \zeta_E=\zeta_{\mathbb Q(\omega)}(\zeta_F/\zeta)^2,
 \qquad |d_E|=27\cdot269^4.
\]
Comparing residues using $w_E=6$, $w_F=2$, and $h_{\mathbb Q(\omega)}=1$ gives
$h_ER_E=h_F^2R_F^2$.
Therefore $h_E=qh_F^2/3$. Integrality and $h_F=7$ force $q=3$ and $h_E=49$.

Finally $\operatorname{Cl}(E)^\sigma=0$.
For a fixed class, norm to $\mathbb Q(\omega)$ followed by extension is both zero and multiplication by three; three is invertible on a group of order $49$.
If $\operatorname{Cl}(E)$ were cyclic, its automorphism group would be abelian.
The relation $c\sigma c=\sigma^{-1}$ would force the action of $\sigma$ to equal its inverse and hence to be trivial, contradicting the absence of fixed classes.
Thus $\operatorname{Cl}(E)\simeq C_7\oplus C_7$.
\end{proof}

% Full proof and printed finite witnesses for the companion supplement.
% Requires amsmath, amssymb, amsthm, booktabs, xurl, hyperref.
\section{Exact stability at seven and all cyclotomic characters}
\label{seven:section}

Fukuda's stability theorem reduces a class-group calculation in an entire
$\mathbb Z_p$-extension to two successive finite layers
\cite[Theorem~1(1), p.~264]{sevenFukuda1994}.
We give its short finite-level proof in the one-prime situation needed here
and certify the two arithmetic inputs. This gives an exact finite Iwasawa
module, including its finite submodule, and hence determines every strict
cyclotomic specialization.

Put
\[
 F=\mathbb Q(\alpha),\quad \alpha^3=269,\qquad
 K=\mathbb Q(\omega),\quad E=KF,
\]
where $\alpha$ is the real root and $\omega$ is a primitive cube root of unity.
Let $Q_n$ be the degree-$7^n$ layer of the cyclotomic
$\mathbb Z_7$-extension of $\mathbb Q$, and set
$F_n=FQ_n$, $K_n=KQ_n$, and $E_n=EQ_n$.
Write $A(k)=\operatorname{Cl}(k)\{7\}$, and
\[
 \Delta=\operatorname{Gal}(E/K)=\langle\delta\rangle,\qquad
 \Gamma=\operatorname{Gal}(E_\infty/E).
\]
Both faithful characters of $\Delta$ take values in $\mu_3\subset\mathbb Z_7$.
Choose $\gamma\in\Gamma$ with
$\langle\varepsilon_7\rangle(\gamma)=8$, and put
$T=\gamma-1$ and $\Lambda=\mathbb Z_7[[T]]$.

For a finite extension $V/\mathbb Q_7$, with integer ring $\mathcal O$, and
a continuous character $\kappa:\Gamma\longrightarrow\mathcal O^\times$,
let
\[
 D_{\chi,\kappa}=(V/\mathcal O)(\chi\kappa).
\]
The characters here are inflated through
$\operatorname{Gal}(E_\infty/K)=\Delta\times\Gamma$.
The strict Selmer group means
\[
 \operatorname{Sel}_{\rm str}(K,D)=
 \ker\left(
 H^1(K,D)\longrightarrow
 \prod_{v\mid7}H^1(K_v,D)
 \times\prod_{v\nmid7}H^1(I_v,D)\right).
 \tag{S1}\label{seven:strict-definition}
\]
The second product is over finite places; imposing the usual condition at
infinite places has no effect on seven-primary cohomology.

\begin{theorem}\label{seven:main}
For every $n\ge0$,
\[
 A(F_n)\simeq\mathbb Z/7\mathbb Z.
 \tag{S2}\label{seven:class-stability}
\]
All norms between these groups are isomorphisms and all cyclotomic actions
are trivial. For either faithful $\chi$ of $\Delta$,
\[
 \left(\varprojlim_n A(E_n)\right)^\chi
 \simeq\Lambda/(7,T).
 \tag{S3}\label{seven:iwasawa}
\]
For every continuous $\kappa$ as above, with $u=\kappa(\gamma)$,
\[
 \operatorname{Sel}_{\rm str}(K,D_{\chi,\kappa})
 \simeq\mathcal O/(7,u-1).
 \tag{S4}\label{seven:all-characters}
\]
These assertions are unconditional.
\end{theorem}

The finite certificate in this proof uses exact algebraic integers, a
matrix of residue characters, and bounded convergent series. A list of
units alone does not certify completeness of a class group or a unit
basis. The certificate proves precisely the seven-part required in
\eqref{seven:class-stability}; it does not need a determination of every
prime divisor of either class number.

\section*{Proof overview}

We first compute the quotient of each class number by the index of an
explicit group of units. A universal regulator lower bound makes the
possible denominator finite. Rigorous intervals isolate the quotient as
the integer seven, while independent residue characters prove that the
unit index is prime to seven. In the cyclotomic tower of $F$, exactly one
prime ramifies. Class field theory therefore identifies class-group
coinvariants with the lower class group. Equality at the first two levels
and Nakayama's lemma force stability at every level. The two faithful
cubic-character parts over $E$ are recovered from the conjugation-fixed
part. Finally, inflation and restriction preserve both strict local
conditions because the faithful cubic character has no invariants.
The resulting specialization of the actual finite module is
\eqref{seven:all-characters}.

\subsection{A finite certification principle}
\label{seven:finite-principle}

\begin{lemma}\label{seven:isolation}
Let $k$ be a number field with regulator $R_k>1/4$. Let $U_0$ contain all
roots of unity of $k$ and be generated modulo torsion by $r$ independent
units, where $r$ is the unit rank. Put
$m=[\mathcal O_k^\times:U_0]$ and $R_0=mR_k$.
Suppose $B$ is an integer with $m<B$ and an interval calculation proves
\[
 \left|\frac{h_k}{m}-7\right|<\frac1B.
 \tag{S5}\label{seven:isolation-bound}
\]
Then $h_k=7m$. If $r$ homomorphisms
$\mathcal O_k^\times\longrightarrow\mathbb F_7$ have a nonsingular
evaluation matrix on the specified units, then $7\nmid m$, and
$A(k)\simeq\mathbb Z/7\mathbb Z$.
\end{lemma}

\begin{proof}
The relation $R_0=mR_k$ follows by comparing the covolumes of the logarithmic
unit lattices. Since $h_k-7m$ is an integer of absolute value less than
$m/B<1$, it vanishes. If $7$ divided $m$, the lattice inclusion
$U_0/\mu(k)\subset\mathcal O_k^\times/\mu(k)$ would have a singular
reduction modulo seven. Equivalently, there would be a product of the
specified units, with exponents not all divisible by seven, which is a
seventh power of a unit. Every displayed homomorphism would kill this
product, contradicting the nonsingular matrix. Thus
$v_7(h_k)=1$, as asserted.
\end{proof}

The regulator bound used here is Friedman's
\cite[Theorem~B, pp.~599--600]{sevenFriedman1989}: $R_k>1/4$ outside three
exceptional totally complex sextic fields. Both fields used below have real
embeddings and degrees three and twenty-one, so this exception is irrelevant.
In particular, one may take $B=\lceil4R_0\rceil$, with an upper enclosure of
$R_0$ used when calculating the ceiling.

For a rational prime $\ell\equiv1\pmod7$ and a prime
$\mathfrak l$ of $k$ with residue field $\mathbb F_\ell$, choose
$z\in\mathbb F_\ell^\times$ of order seven. The character
\[
 u^{(\ell-1)/7}\equiv z^{\,\ell_{\mathfrak l}(u)}
 \pmod{\mathfrak l},\qquad
 \ell_{\mathfrak l}(u)\in\mathbb F_7
 \tag{S6}\label{seven:residue-character}
\]
is a homomorphism on all units. This supplies the homomorphisms in the lemma.
The fields here have a real embedding and odd degree, so their roots of
unity are exactly $\{\pm1\}$.

\subsection{The zeta residues as bounded finite sums}
\label{seven:analytic}

Set $N=217083=3\cdot269^2$, and let $\rho$ be the standard two-dimensional
representation of $\operatorname{Gal}(E/\mathbb Q)\simeq S_3$. Then
\[
 L_0(s)=L(s,\rho)=\frac{\zeta_F(s)}{\zeta(s)},\qquad
 D_F=-N.
 \tag{S7}\label{seven:base-L}
\]
For completeness, the discriminant exponent at $269$ is two by tame total
ramification. At three, $v_3((-1)^3-269)=3>2v_3(3)=2$, so Hensel's
lemma gives a root in $\mathbb Q_3$ and
$F\otimes\mathbb Q_3=\mathbb Q_3\times\mathbb Q_3(\omega)$.
The exponent at three is therefore one. These observations also identify
the ramified Euler factors below.

Let $\psi$ be the order-seven character modulo $49$ with
$\psi(3)=\zeta_7$, and put $L_j=L(s,\rho\otimes\psi^j)$.
For $1\le j\le6$, the character $\psi^j$ is even and primitive of conductor
$49$: $3$ has order $42$ modulo $49$, and $-1=3^{21}$.
Its conductor cannot divide seven because
$|(\mathbb Z/7\mathbb Z)^\times|=6$.
Since $\rho$ is unramified at seven, twisting its two unramified dimensions
gives conductor exponent four there. Thus
\[
 N_0=N,\quad N_j=49^2N\ (j\ne0),\qquad
 \Lambda_j(s)=N_j^{s/2}\Gamma_{\mathbb C}(s)L_j(s),
 \quad
 \Gamma_{\mathbb C}(s)=2(2\pi)^{-s}\Gamma(s).
 \tag{S8}\label{seven:conductors}
\]
The Dedekind functional equation gives
$\Lambda_0(s)=\Lambda_0(1-s)$
\cite[Remark~18.14]{sevenSutherland2015}.
The representation is induced from a nontrivial cubic Hecke character of
$K$ of conductor $(269)$, whose conjugate is its distinct inverse.
Hecke's theta construction therefore gives the cuspidal weight-one form
$f(z)=\sum_{n\ge1}a_ne^{2\pi inz}$ of level $N$ and nebentypus
$\eta=\left(\frac{-3}{\,\cdot\,}\right)$; see
\cite[Theorem~5]{sevenSimard2015} for the theta transformation and
\cite[\S3]{sevenDarmonLauderRotger2015} for the weight-one induced
representation and cuspidality criterion.

Here is a derivation of the twisting sign, to fix its orientation.
For a primitive even character $\xi\bmod M$, $M=49$, set
$\tau(\xi)=\sum_{a\bmod M}\xi(a)e^{2\pi ia/M}$ and
\[
 A_a=\begin{pmatrix}1&a/M\\0&1\end{pmatrix},\qquad
 W_Q=\begin{pmatrix}0&-1\\Q&0\end{pmatrix}.
\]
With the normalized weight-one slash operator, finite Fourier inversion gives
\[ f_\xi=\tau(\bar\xi)^{-1}\sum_a\bar\xi(a)f|A_a. \]
Choose $b$ with $Nab\equiv-1\pmod M$. Direct multiplication gives
\[
 A_aW_{NM^2}=
 \begin{pmatrix}(1+Nab)/M&a\\Nb&M\end{pmatrix}M W_NA_b.
 \tag{S9}\label{seven:fricke-matrix}
\]
The first matrix has determinant one, lies in $\Gamma_0(N)$, and has
lower-right entry $M$, with $\eta(M)=1$.
Also $\bar\xi(a)=\xi(N)\xi(b)$.
The base functional equation gives $f|W_N=-if$ by Mellin inversion.
Using $\tau(\xi)\tau(\bar\xi)=M$ consequently gives
\[
 \Lambda_j(s)=\varepsilon_j\Lambda_{7-j}(1-s),\qquad
 \varepsilon_j=\psi^j(N)\frac{\tau(\psi^j)^2}{49}.
 \tag{S10}\label{seven:epsilon}
\]
There is a short exact verification of the Gauss sums.
The relation $3^6=1-7\bmod49$ gives
$\psi(1+7b)=\zeta_7^b$. Grouping terms as $a+7t$ shows that only
$a\equiv-j\pmod7$ survives in $\tau(\psi^j)$.
Since $N\equiv3^{33}\pmod{49}$, the results are
\[
\begin{array}{c|rrrrrr}
 j&1&2&3&4&5&6\\ \hline
 \tau(\psi^j)/7=\zeta_{49}^{\,b_j}
   &48&19&18&31&30&1\\
 \varepsilon_j=\zeta_{49}^{\,c_j}
   &33&10&43&6&39&16
\end{array}
 \tag{S11}\label{seven:gauss-table}
\]
where the entries in the last two rows are the exponents $b_j,c_j$.

The coefficients $a_n$ are multiplicative and have the following complete
prime-power specification, including $a_{p^0}=1$:
\[
\begin{array}{c|l}
 p& a_{p^k}\quad(k\ge1)\\ \hline
269&0\\
3&1\\
p\equiv2\pmod3,\ p\ne269
   &1\text{ for even }k,\quad0\text{ for odd }k\\
p\equiv1\pmod3,\ 269^{(p-1)/3}\equiv1\pmod p&k+1\\
p\equiv1\pmod3,\ 269^{(p-1)/3}\not\equiv1\pmod p
   &1,-1,0\text{ for }k\equiv0,1,2\pmod3 .
\end{array}
 \tag{S12}\label{seven:coefficients}
\]
These are the Euler factors for the identity, transposition, and
three-cycle classes, together with the inertia invariants at the two
ramified primes. In particular, $|a_n|\le d(n)\le n$.

Put $\lambda_j=2\pi/\sqrt{N_j}$.
Splitting the Mellin integral at one and using
\eqref{seven:epsilon} yields, with
$E_1(x)=\int_x^\infty e^{-t}\,dt/t$,
\[
 L_j(1)=
 \sum_{n\ge1}\frac{a_n\psi^j(n)e^{-\lambda_jn}}{n}
 +\varepsilon_j \lambda_j
 \sum_{n\ge1}a_n\overline{\psi^j(n)}E_1(\lambda_jn).
 \tag{S13}\label{seven:L1-formula}
\]
For $j=0$ the character in this formula is identically one, including at
multiples of seven, and $\varepsilon_0=1$.
Indeed, the theta transformation is
$\Theta_j(t)=\varepsilon_jt^{-1}\Theta_{7-j}(1/t)$; after changing
$t$ to $1/t$ in the lower Mellin integral, its value at $s=1$ is precisely
the second term in \eqref{seven:L1-formula}.
Because $0<E_1(x)\le e^{-x}/x$, deleting all terms $n>X$ in both sums
has joint absolute error at most
\[
 {\cal T}(c,X)=\frac{2e^{-c(X+1)}}{1-e^{-c}}.
 \tag{S14}\label{seven:tail}
\]
This estimate uses only $|a_n|\le n$ and a geometric series.

The degree-seven field $Q_1$ has
$\zeta_{Q_1}(s)=\prod_{j=0}^6L(s,\psi^j)$.
Here $\psi^0$ means the trivial character of conductor one, whose factor is
$\zeta(s)$, so its Euler factor at seven is retained.
For $j\ne0$ its factors at one are given by the finite identity
\[
 L(1,\psi^j)=
 -\frac{\tau(\psi^j)}{49}
 \sum_{\substack{1\le a<49\\7\nmid a}}
       \overline{\psi^j(a)}
       \log\left(2\sin\frac{\pi a}{49}\right).
 \tag{S15}\label{seven:dirichlet-residue}
\]
This follows by Fourier expansion of $-\log(2\sin\pi x)$ and the
Gauss-sum identity; evenness combines the two Fourier directions.
Pairing conjugate characters makes the resulting residue real and positive.
Finally, the permutation representation for $F_1/\mathbb Q$ is
$(1+\rho)\otimes\bigoplus_{j=0}^6\psi^j$, so Artin formalism gives
\[
 \operatorname*{res}_{s=1}\zeta_{F_1}(s)
 =\operatorname*{res}_{s=1}\zeta_{Q_1}(s)
       \prod_{j=0}^6L_j(1).
 \tag{S16}\label{seven:residue-factorization}
\]
This identity includes every Euler factor, including those at seven.

\subsection{Explicit units and rigorous interval isolation}
\label{seven:certification}

Let
\[
 q(Y)=Y^7-21Y^5-21Y^4+91Y^3+112Y^2-84Y-97,
 \qquad
 \beta=\sum_{a\in\{1,18,19,30,31,48\}}\zeta_{49}^{a}.
 \tag{S17}\label{seven:period}
\]
The set of exponents is the subgroup of order six in
$(\mathbb Z/49\mathbb Z)^\times$. Thus $\mathbb Q(\beta)=Q_1$.
The polynomial $q$ is irreducible and $q(\beta)=0$ by direct arithmetic
modulo $\Phi_{49}$.
Write $\theta=\alpha-\beta$. The polynomial $p$ of degree twenty-one and
the rational function $P/Q$ printed in the finite data below satisfy
\[
 p(\theta)=0,\qquad \alpha=P(\theta)/Q(\theta),\qquad
 \alpha^3=269,\qquad q(\alpha-\theta)=0.
 \tag{S18}\label{seven:field-identity}
\]
Irreducibility of $p$ and $q$, or the relatively prime degrees three and
seven together with these identities, identifies this field as $F_1$.
The automorphism $\tau$ fixes $\alpha$ and is
$\zeta_{49}\mapsto\zeta_{49}^3$ on $Q_1$; its action on $\beta$ is the
degree-six rational polynomial printed below.

The first six units $u_1,\ldots,u_6$ lie in $\mathbb Q(\beta)$ and are
given by their seven coefficients and denominators.
The seventh is
\[
 u_7=-\varepsilon,\qquad
 \varepsilon=270-739\alpha+108\alpha^2.
 \tag{S19}\label{seven:base-unit}
\]
For the explicitly printed $\eta$, the remaining units are
$u_{8+i}=\tau^i\eta$ for $0\le i\le5$.
Exact rational polynomial arithmetic verifies that all thirteen units are
integral and have absolute norm one. One can check integrality either in
the certified maximal order or from the monic integral characteristic
polynomial of multiplication by the element. The discriminant and
signature are
\[
 D_{F_1}=-N^7\,7^{36},\qquad (r_1,r_2)=(7,7).
 \tag{S20}\label{seven:disc-signature}
\]
Independently, $D_{Q_1}=49^6=7^{12}$ by the conductor--discriminant
formula. Its discriminant is coprime to $D_F$, so the product of the
integral bases is a maximal order in the compositum and gives
\eqref{seven:disc-signature}.

Here are complete numerical specifications for the interval calculation.
Generate the integer coefficients in \eqref{seven:coefficients} by a
prime-factor sieve. In \eqref{seven:L1-formula} take $X=14000$ for $j=0$
and $X=550000$ for $j\ne0$, and evaluate all terms by real and complex
ball arithmetic at precision $300$ bits.
The bounds \eqref{seven:tail} are, respectively,
$1.496\cdot10^{-80}$ and $1.328\cdot10^{-62}$.
The resulting enclosures, rounded outward, include
\[
\begin{split}
 L_0(1)&\in
 [\,1.6976757232625543059851766859761711810674
       \ \pm 10^{-40}\,],\\
 \prod_{j=0}^6L_j(1)&\in
 [\,0.6494049966247962459428741324547233673973
       \ \pm 10^{-40}\,].
\end{split}
 \tag{S21}\label{seven:analytic-intervals}
\]
For the regulator, isolate the seven disjoint real roots of $q$ with
complex ball arithmetic, take
$\alpha_0=\exp(\log269/3)$ and
$\alpha_1=\alpha_0(-1+i\sqrt3)/2$, and evaluate the thirteen units at
$(\alpha_0,\beta_j)$ and $(\alpha_1,\beta_j)$.
Form the logarithmic matrix with entries $\log|\sigma(u_i)|$ for real
embeddings and $2\log|\sigma(u_i)|$ for complex pairs, delete one real
row, and take the absolute determinant.
At $2048$ bits this gives the regulator $R_0$ of the displayed unit group
with over $600$ bits of relative accuracy. Arb's enclosures, including
the root isolation and determinant operations, are used throughout
\cite{sevenJohansson2017}.
Rounded outward at lower precision, the relevant values are
\[
\begin{split}
 R_0&\in
 [\,8.7214139234209020456623583408322697739521
       \ \pm10^{-40}\,]\,10^{24},\\
 \operatorname*{res}_{s=1}\zeta_{Q_1}(s)
   &\in[\,0.2996769180948235445839119514335675929714
       \ \pm10^{-40}\,].
\end{split}
 \tag{S22}\label{seven:regulator-intervals}
\]
The positive determinant proves independence of the units.

The analytic class-number formula, in the convention just used, is
\[
 \operatorname*{res}_{s=1}\zeta_k(s)
  =\frac{2^{r_1}(2\pi)^{r_2}h_kR_k}
         {|\mu(k)|\sqrt{|D_k|}}
 \tag{S23}\label{seven:class-formula}
\]
\cite[Theorem~18.13]{sevenSutherland2015}.
If $m_1$ is the index of the displayed thirteen units and $\{\pm1\}$,
\eqref{seven:residue-factorization} and the unrounded ball evaluations give
\[
 \frac{h_{F_1}}{m_1}
 =\frac{\sqrt{N^7\,7^{36}}}
        {2^{13}\pi^7R_0}
   \operatorname*{res}_{s=1}\zeta_{Q_1}(s)\prod_{j=0}^6L_j(1)
 \in[\,7\pm1.11\cdot10^{-59}\,].
 \tag{S24}\label{seven:quotient-interval}
\]
The stated precision and cutoffs, rather than substitution of the rounded
displayed decimals, specify this computation.
Friedman's bound gives
\[
 m_1<B=34885655693683608182649434,
 \qquad
 B\left|h_{F_1}/m_1-7\right|<3.88\cdot10^{-34}<1.
 \tag{S25}\label{seven:index-bound}
\]
The thirteen residue primes and their full evaluation matrix are printed
below. Their determinant is $1$ in $\mathbb F_7$.
Lemma~\ref{seven:isolation} proves $v_7(h_{F_1})=1$.

The same short certificate supplies the base layer without assuming a
class-number search. For \eqref{seven:base-unit}, exact arithmetic gives
$N_{F/\mathbb Q}(\varepsilon)=1$ and
\[
 0<R_\varepsilon=|\log|\varepsilon||<18,\qquad
 \frac{\sqrt N L_0(1)}{2\pi R_\varepsilon}
 \in[\,7\pm4.0\cdot10^{-68}\,].
 \tag{S26}\label{seven:base-quotient}
\]
Thus $m_0=[\mathcal O_F^\times:\langle-1,\varepsilon\rangle]<72$ and
$h_F=7m_0$.
At the prime $491$ the map $\alpha\mapsto444$ gives
$\varepsilon\mapsto88$, while
\[
 138^7\equiv1,\quad138\not\equiv1,\qquad
 88^{70}\equiv223\equiv138^4\not\equiv1\pmod{491}.
 \tag{S27}\label{seven:base-saturation}
\]
This nonzero residue character proves $7\nmid m_0$.
Consequently
\[
 A(F)=A(F_1)=\mathbb Z/7\mathbb Z.
 \tag{S28}\label{seven:finite-inputs}
\]
All integers, unit expressions, coefficient rules, cutoffs and tail bounds
needed for these two assertions have now been specified.

\subsection{Norm control and stabilization}
\label{seven:stability-proof}

The cube classes in $\mathbb F_7$ are $0,1,6$, whereas $269\equiv3$.
Thus $X^3-269$ is irreducible modulo seven, and seven is unramified and inert
in $F$. The local extension at this unique prime is the unramified cubic
extension of $\mathbb Q_7$.
Each cyclotomic extension $Q_n/Q_m$ is totally ramified at seven and
unramified everywhere else. Base change by the unramified cubic local
field preserves this total ramification and its degree. Therefore every
$F_n/F_m$ has exactly one ramified finite prime, which is totally ramified.

\begin{lemma}\label{seven:norm-control}
Let $L/k$ be cyclic of degree $7^r$, with exactly one ramified finite prime,
which is totally ramified. Then the ideal norm induces an isomorphism
\[
 A(L)_{\operatorname{Gal}(L/k)}\ \xrightarrow{\ \sim\ }\ A(k).
 \tag{S29}\label{seven:coinvariant-control}
\]
\end{lemma}

\begin{proof}
Let $H$ be the Hilbert seven-class field of $L$. It is Galois over $k$,
since its defining maximality property is invariant under every
automorphism of $L/k$. Set
$G=\operatorname{Gal}(H/k)$ and
$A=\operatorname{Gal}(H/L)=A(L)$.
An inertia group $I$ above the sole ramified prime has order $7^r$:
$H/L$ is unramified and $L/k$ is totally ramified.
It therefore maps isomorphically to $\operatorname{Gal}(L/k)$, giving
$G=A\rtimes I$.
If $\sigma$ generates $I$, then $[G,G]=(\sigma-1)A$.
Killing $I$ and $[G,G]$ kills all inertia groups and gives the maximal
abelian unramified quotient of $G$, namely $A/(\sigma-1)A$.
The Hilbert seven-class field $H_k$ of $k$ is contained in $H$, since
$H_kL/L$ is unramified, abelian and of seven-power degree. Conversely,
every abelian unramified subextension of $H/k$ lies in $H_k$.
The quotient is therefore $A(k)$.
Under reciprocity this restriction map is the ideal norm, proving the
claim. There is no archimedean inertia of odd order to consider.
\end{proof}

This is the one-prime case of the usual class-field-theoretic norm
calculation; compare
\cite[Corollaries~1.3.9 and 1.3.12]{sevenSharifi}.
The proof includes the finite control map, so that no unidentified
finite kernel is carried into the next argument.

Fix $n\ge1$, put $A_n=A(F_n)$, and regard it as a module over
$R_n=\mathbb Z_7[C_{7^n}]$.
This ring is local with maximal ideal $(7,T)$, because
$R_n/7R_n=\mathbb F_7[T]/(T^{7^n})$.
Lemma~\ref{seven:norm-control} gives
\[
 A_n/TA_n\simeq A(F),\qquad
 A_n/\omega_1(T)A_n\simeq A(F_1),\qquad
 \omega_1(T)=(1+T)^7-1=T\nu(T),
 \tag{S30}\label{seven:two-quotients}
\]
where
$\nu(T)=7+21T+35T^2+35T^3+21T^4+7T^5+T^6$.
For $n=1$ the second quotient is simply $A_1$, so the statement remains
valid. Both quotients have order seven by \eqref{seven:finite-inputs}.
Since $\omega_1A_n\subseteq TA_n$, equality of their indices implies
\[
 TA_n=\omega_1A_n=\nu(T)\,TA_n.
\]
But $\nu(T)\in(7,T)$, so Nakayama's lemma applied to the finite
$R_n$-module $TA_n$ gives $TA_n=0$.
The first quotient in \eqref{seven:two-quotients} now proves
\eqref{seven:class-stability}.
Applying Lemma~\ref{seven:norm-control} to every $F_n/F_m$ shows that all
norms are surjective, hence isomorphisms between groups of order seven.
The resulting inverse limit is $\Lambda/(7,T)$.
This also proves directly the needed case of Fukuda's theorem
\cite[Theorem~1(1)]{sevenFukuda1994}.

The same norm lemma applied to $Q_n/\mathbb Q$, followed by Nakayama,
gives
\[
 A(Q_n)=0 \qquad(n\ge0).
 \tag{S31}\label{seven:rational-layer}
\]
Indeed its coinvariants are $A(\mathbb Q)=0$.

\subsection{The two faithful cubic-character components}
\label{seven:character-components}

Let $c$ be complex conjugation. Then
$\operatorname{Gal}(E_n/Q_n)=S_3$ with
$c\delta c=\delta^{-1}$, and $E_n^{\langle c\rangle}=F_n$.
Set $C_n=A(E_n)$.
For an extension of degree $d$ prime to seven, extension and norm of ideals
satisfy
$Nj=d$ and $jN=\sum_g g$.
Since $d$ acts invertibly on seven-primary groups, extension identifies the
lower class group with the invariant subgroup, without a ramification
restriction. Thus
\[
 C_n^{c=1}\simeq A(F_n),\qquad
 C_n^\Delta\simeq A(K_n),\qquad
 (C_n^\Delta)^{c=1}\simeq A(Q_n)=0.
 \tag{S32}\label{seven:prime-to-seven-descent}
\]
For $\varphi\in\{1,\chi,\chi^{-1}\}$ put
$e_\varphi=\frac13\sum_{d\in\Delta}\varphi(d)^{-1}d$.
The idempotents split
$C_n=C_n^1\oplus C_n^\chi\oplus C_n^{\chi^{-1}}$, and $c$ exchanges the
last two summands.
Since the $c$-fixed part of $C_n^1$ vanishes, projection gives an
isomorphism
\[
 e_\chi:C_n^{c=1}\ \xrightarrow{\ \sim\ }\ C_n^\chi,
 \qquad x\longmapsto e_\chi x,
 \qquad y\longmapsto y+cy
 \tag{S33}\label{seven:projection}
\]
with the displayed inverse.
Extension of ideals, the idempotents, and these projections commute with
the cyclotomic action and with the norm maps in the compositum tower.
Equations \eqref{seven:class-stability} and \eqref{seven:projection}
therefore prove \eqref{seven:iwasawa} for both faithful characters.

Only $(C_n^1)^{c=1}=0$ was used; deleting $C_n^1$ itself would be false.
For example, $K$ has class number one by its Minkowski bound
$(2/\pi)\sqrt3<2$. In $K_n/K$ the two primes above seven are totally
ramified, and each unit of $K$ is the norm of a unit, since the norm on
$\mu_6$ is raising to $7^n$.
The cyclic ambiguous class-number formula consequently gives
$|A(K_n)^{\operatorname{Gal}(K_n/K)}|=7^n$.
This growth lies in the conjugation-odd part and does not affect
\eqref{seven:projection}; the relevant cyclic class-field exact sequence
is given in \cite[Theorem~1.3.16]{sevenSharifi}.

\subsection{Strict local control and specialization}
\label{seven:control}

Seven splits into two primes in $K$, and $E/K$ is unramified and inert of
degree three at each of them. Hence each prime over seven in $E_n$ is fixed
by $\Delta$.
If
\[
 C_n'=\operatorname{Cl}(\mathcal O_{E_n}[1/7])\{7\},
 \qquad M'=\varprojlim_n C_n',
\]
the kernel of $C_n\to C_n'$ is generated by the seven-prime classes and
is $\Delta$-trivial. It follows that
\[
 (M')^\chi
 \simeq\left(\varprojlim_n C_n\right)^\chi
 \simeq\Lambda/(7,T).
 \tag{S34}\label{seven:split-module}
\]
Class field theory identifies $M'$ with the maximal unramified abelian
pro-seven Galois group over $E_\infty$ in which every seven-prime splits
completely. This is the norm inverse limit of the finite Hilbert
$S$-class-field correspondences
\cite[Remark~1.3.21 and Proposition~3.2.20]{sevenSharifi}.

Here one can check the global and local descent maps explicitly.
Let $G=\Delta\times\Gamma$ and $D=D_{\chi,\kappa}$.
The element $\chi(\delta)-1$ is a unit in $\mathcal O$, so $D^\Delta=0$.
Cohomology in positive degree for the order-three group $\Delta$ vanishes
on seven-primary modules, by averaging.
The Hochschild--Serre sequence therefore gives
$H^i(G,D)=0$ for all $i\ge0$.
The inflation--restriction sequence is consequently an isomorphism
\[
 H^1(K,D)\ \xrightarrow{\ \sim\ }\
 H^1(E_\infty,D)^G
 =\operatorname{Hom}_{\rm cts}(G_{E_\infty},D)^G.
 \tag{S35}\label{seven:global-inflation}
\]
At either seven-prime the decomposition group in $G$ is all of
$\Delta\times\Gamma$: the local extension is the compositum of the
unramified cubic extension and the totally ramified cyclotomic tower.
The identical argument shows that local inflation--restriction is an
isomorphism there. Thus a class vanishes on the full local decomposition
group before descent if and only if its descended homomorphism does.
At a finite prime away from seven the inertia quotient in $G$ has order
one or three; its positive cohomology again vanishes. This makes
restriction injective on the inertia cohomology, so the unramified
condition is also equivalent before and after descent.
These are precisely the two local conditions in
\eqref{seven:strict-definition}. We obtain
\[
 \operatorname{Sel}_{\rm str}(K,D)
 =\operatorname{Hom}_{\rm cts}(M',D)^{\Delta\times\Gamma}.
 \tag{S36}\label{seven:Hom-description}
\]

For clarity about character and inverse conventions, equivariance means
\[
 f(\delta x)=\chi(\delta)f(x),\qquad f(\gamma x)=u f(x).
 \tag{S37}\label{seven:orientation}
\]
The first condition selects the natural $\chi$-component of $M'$.
On that component $\gamma=1$ and $7=0$ by \eqref{seven:split-module}.
Hence \eqref{seven:Hom-description} is $D[7,u-1]$, which, since
$\mathcal O$ is a discrete valuation ring, is noncanonically
$\mathcal O/(7,u-1)$.
This proves \eqref{seven:all-characters}.
Writing a contragredient convention instead replaces $u-1$ by
$u^{-1}-1=-u^{-1}(u-1)$, the same ideal. No Tate twist or unspecified
specialization kernel occurs in this argument.

\subsection{Weights, wild characters, and coefficient normalization}
\label{seven:weights}

For $e\in\mathcal O$, the convergent exponential gives the continuous
character determined by
$\kappa(\gamma)=\exp(-e\log8)$.
It satisfies $u-1\in7\mathcal O$. In particular this includes
$\kappa=\langle\varepsilon_7\rangle^{-e}$ for $e\in\mathbb Z_7$, including
$e=0$, and
\[
 \operatorname{Sel}_{\rm str}
   \bigl(K,(V/\mathcal O)(\chi\langle\varepsilon_7\rangle^{-e})\bigr)
 \simeq\mathcal O/7\mathcal O.
 \tag{S38}\label{seven:principal-weights}
\]
Over $\mathbb Z_7$ every such group has length one.
Over a ramified coefficient ring its $\mathcal O$-length is
$e(V/\mathbb Q_7)$; it must not be replaced by a length-one quotient by
a uniformizer.

More generally, suppose $\mathcal O$ contains a primitive $7^r$th root
$\zeta$, where $r\ge1$, and put
$u=\zeta\exp(-e\log8)$ with $e\in\mathcal O$.
Since
$v_7(\zeta-1)=1/(6\cdot7^{r-1})<1$
and $\exp(-e\log8)-1\in7\mathcal O$, the two summands in
$u-1=(\zeta-1)+\zeta(\exp(-e\log8)-1)$ have different valuations.
Thus
\[
 \operatorname{Sel}_{\rm str}(K,D_{\chi,\kappa})
 \simeq\mathcal O/(\zeta-1),\qquad
 \frac{\log\#\operatorname{Sel}_{\rm str}(K,D_{\chi,\kappa})}
      {[V:\mathbb Q_7]}
 =\frac{\log7}{6\cdot7^{r-1}}.
 \tag{S39}\label{seven:wild-cost}
\]
The valuation formula follows by applying the Eisenstein polynomial
$\Phi_{7^r}(1+X)$ to $\zeta-1$.
For the principal branch \eqref{seven:principal-weights}, the same
normalized logarithmic size is $\log7$.
For an arbitrary continuous $\kappa$, the completely general answer
remains \eqref{seven:all-characters}; no parametrization of all such
characters by $\exp(-e\log8)$ is assumed.

Finally consider a coefficient family written in the original finite
character convention as
\[
 D_{j,a,e}=(V/\mathcal O)
    \bigl(\chi^j\xi_7^a\varepsilon_7^{-e}\bigr),\qquad
 j\in\{1,2\},\quad0\le a<6,\quad e\ge1,
 \qquad \xi_7=\omega_7^v,\quad v\in\{1,5\}.
 \tag{S40}\label{seven:matched-coefficients}
\]
If the finite residual characters are matched, then
$\xi_7^a\omega_7^{-e}=1$ exactly, because reduction is injective on
$\mu_6$. Equivalently, $va\equiv e\pmod6$, which determines one $a$
for every $e$. Thus each matched coefficient is exactly
$(V/\mathcal O)(\chi^j\langle\varepsilon_7\rangle^{-e})$, and
\[
 s(e)=1\quad\hbox{for every positive matched weight},\qquad
 \sum_{j=1}^{2}\sum_{e=1}^{N}
 \frac{\log\#\operatorname{Sel}_{\rm str}(K,D_{j,a(e),e})}
      {[V:\mathbb Q_7]}
 =2N\log7.
 \tag{S41}\label{seven:matched-cost}
\]
There is no exceptional weight in this matched family. Statements about
other tame characters or additional factors require their own local
matching and are not used here.

\subsection{Complete finite unit and residue data}
\label{seven:finite-data}

The following arrays are exact integers. Coefficients for $p$, $q$, $P$ and
$Q$ are in descending order. The coefficients for $\tau$, the six small
units, and the cubic unit are in ascending order. A small-unit numerator
$[a_0,\ldots,a_6]$ represents $\sum a_i\beta^i$; the cubic-unit numerator
represents a polynomial in $\alpha$.
The twenty-one coefficients of $\eta$ are ordered with $\alpha$ varying
first: coefficient number $3j+i+1$ multiplies $\alpha^i\beta^j$,
$0\le i\le2$, $0\le j\le6$.
To keep every integer on the page, each coefficient of $\eta$ is printed as
$[s,a_0,a_1,a_2]$ in base $B_*=10^{100}$, meaning
$s(a_0+B_*a_1+B_*^2a_2)$; no leading zeroes are required.

Each residue triple is $(\ell,r,z)$: reduce $\theta$ to $r$ modulo $\ell$
and use $z$ in \eqref{seven:residue-character}. The thirteen primes with
these residues are distinct, even when their rational primes coincide.
The roots of $p$ are simple at all these residues; the derivative residues,
in the same order, are
\[
 153,14,99,174,74,74,436,46,119,368,219,76,606.
\]
Hensel's lemma therefore supplies the corresponding maps of the full
integer ring to $\mathbb F_\ell$.
The rows of the final matrix are in this order and the columns are
$u_1,\ldots,u_{13}$.
Polynomial substitution, the norm and integrality tests, the finite-field
powers in \eqref{seven:residue-character}, and Gaussian elimination modulo
seven are a complete exact verification of these finite witnesses.

\begingroup
\fontsize{7.5}{9.5}\selectfont
\begin{verbatim}
p = [1, 0, -63, -1820, 1596, 47859, 1550766, -771102, 74382, -773205433, 265573350, -12672069060,
  225337203987, -164087914221, 5350549479537, -37293576067279, 55623366413664, -844644350170257,
  3089185218698537, -5358907459195815, 42933560367107697, -74554880595863287]
q = [1, 0, -21, -21, 91, 112, -84, -97]
P = [7, 0, -252, 43540, 3115, -767592, 28694829, 4854766, -154925106, 3054427187, 508674397, 875124348,
  32367791232, 6622577365]
Q = [28, 0, -819, 72163, 8253, -926394, 25924017, 4555698, -71871954, 1423159430, 242638326,
  1061561823, 4574293518]
tau_numerator = [3210, -2025, -2700, 820, 495, -51, -15]
tau_denominator = 589
small_unit_numerators = [
  [48, 18, -65, -13, 17, 2, -1],
  [-842, 2689, 5156, -305, -1495, -78, 81],
  [-5185, -1631, 4697, 659, -1028, -69, 49],
  [5927, -5217, -7545, 1048, 1825, 44, -91],
  [5085, -3117, -2389, 743, 330, -34, -10],
  [10258, -972, -10131, -431, 2358, 164, -125],
]
small_unit_denominators = [1, 589, 589, 589, 589, 589]
cubic_unit_numerator = [-270, 739, -108]
eta_denominator = 1767
eta_coefficients_in_base_10_power_100 = [
  [1,
   8110910039399052799388839509913995719324441857259887672553997475894562301601700471535299521506405517,
   3431137402558321943275234065603513747732585590776418713813607417676553074824870235139769709481053259,
   3940696968887],
  [1,
   2863301462649584998067112611608264044687998196095409917598114508277713738730199925757154728888953055,
   970812101337790668521585581275454765494152664299295398136805535577832483461611290913002007827311071,
   610457752147],
  [1,
   4821724478914456225350364309117571242884298504341022610343187741267858848133670184073817508713496298,
   3511990174213849653115720097913645282815157493086558732902294533502899394979501897922579875243606017,
   94566689623],
  [1,
   760344159491231074673603899303458724681328958152289201635437873131003916558596294478134311589863496,
   2317480636152653863897711154536450131210677689030596417135280558908466878562464587662862833652562118,
   7059069004615],
  [1,
   5713646533666935748952660526437906230094439209439484821582261626527132028034044010356200033729959598,
   8155918751592805126188738962782659016016516467436799064429460997991815680374480211693583033342107571,
   1093528233921],
  [1,
   5458376416976996449886152263319535978441113085075418487179464286576609680040631434748810087394632375,
   2815314908475193445673806280168229698608681052505411231685686073970229478053287079417408542790496382,
   169399675453],
  [1,
   9104803622062202912042424298875099840356676073442457598635399038039260264973760983790173980634551253,
   5209978760288616211810004718597728762354940779407360157403598658202633153414464523403435572630242052,
   1981995488635],
  [1,
   3229190277068481099567369991097155293895283837272933560810654867460601672084248677453637066047248351,
   1066602590408197039412448163000624315584233985107205159080701191411246693038208294958734638546267908,
   307033126452],
  [1,
   4719833413501908223670674319560071208620416677277259622958912093004163789224972446770714978198517369,
   3589225343881354158630968166458173644607882286774087545514248915094188084702046502205303007048245172,
   47562843245],
  [-1,
   7252159511744921579747663915866179139747715077640636824080064020836768547568447722702965774255379864,
   2117795437844894019014366078716012207469677439109909489100232555667366018470798129278480577073347195,
   1862911672364],
  [-1,
   4596952397808038061504443355452507844075138976341878918188397033987218237849921882684029382396888813,
   5548014240268422805702685958606065400771312345386887361270938503247549652804788711898670621959818004,
   288585719972],
  [-1,
   8576573868236387713297262546743178973121352627702349510026823016588229660225220810296942340315536894,
   7246874522795187239999206975901393322813195843039386133915845845849949086654946548202311744607548915,
   44705134981],
  [-1,
   7281143800513613602674102889822350672969769041691650309307424554802016538032032999848411121384589417,
   7247656128371841488568688653380385878443917715304128569138890771204319093551264064535444911867233869,
   870696576112],
  [-1,
   2603158096105930003776498071993161225442241308288062551734654529270377694879518302937682334035359418,
   8801224015145560570571877684554110083953909461438848969946722322537320157416843524148670087474475623,
   134880575403],
  [-1,
   2144340361888752919586412430689254661717183112022793764572101308645944000985194209283484734442009651,
   4344877732372528348076072331962740358969404598366277209320255591698654981676044669701093523564674725,
   20894500013],
  [1,
   1658789263126353124561691351357597966888983262121120607405319802389834239184987248275993385042870297,
   4563350942981641714892768486786517151934762869785113709687288930688939203417092386426646047883401037,
   47445317587],
  [1,
   3334821144767026713097794436460203379080665880747592876486755438276015329827178694499770466419371729,
   4734315972881480902169887703698274649856456832930081114001007355628760067288435055252831111892960802,
   7349806938],
  [1,
   2179355790946349106219951244546405670330306327971639576268515952647904221008259146456375776668744595,
   7511295224937829629681218123710093110445808535628505578790220369499924491111462384032633379660007227,
   1138566770],
  [1,
   1940103518950861831403240058854220493243755110925587464386601540881397210941906392843448161215669392,
   7398981907903882144917187999221062721013600253004142759479698501891917306885982478145740682084331734,
   43903316832],
  [1,
   3471460889921431092203467331457652901571380143073529228472752287967796312304491085909952970479760717,
   4781793234571883106444295805054726784965688101009565012732394759877052590133372264885207805147053091,
   6801111660],
  [1,
   1294025870817364042382316796528664510562816139774755081398276045466303764150138262507798612436708108,
   7200960424005516814554303735541605542146140398242153380225569364053345312160643463981887310532844365,
   1053567774],
]
prime_root_zeta_triples = [
  [197, 164, 104],
  [197, 119, 104],
  [197, 110, 104],
  [197, 69, 104],
  [197, 50, 104],
  [197, 29, 104],
  [491, 435, 138],
  [491, 431, 138],
  [491, 411, 138],
  [491, 107, 138],
  [491, 98, 138],
  [491, 93, 138],
  [883, 874, 707],
]
residue_character_matrix = [
  [2, 3, 4, 2, 6, 5, 0, 5, 5, 2, 3, 3, 5],
  [5, 6, 2, 5, 1, 1, 0, 5, 5, 5, 2, 3, 3],
  [0, 1, 3, 1, 0, 2, 0, 5, 5, 5, 5, 2, 3],
  [6, 0, 6, 2, 5, 5, 0, 3, 5, 5, 5, 5, 2],
  [5, 4, 0, 0, 2, 6, 0, 2, 3, 3, 5, 5, 5],
  [2, 5, 1, 5, 4, 0, 0, 3, 3, 5, 5, 5, 5],
  [6, 6, 2, 0, 4, 5, 4, 3, 4, 1, 0, 6, 4],
  [4, 4, 2, 5, 0, 3, 4, 3, 3, 4, 1, 0, 6],
  [5, 2, 0, 3, 2, 0, 4, 1, 0, 6, 4, 3, 3],
  [3, 0, 6, 3, 0, 6, 4, 4, 3, 3, 4, 1, 0],
  [3, 2, 0, 0, 6, 0, 4, 4, 1, 0, 6, 4, 3],
  [0, 0, 4, 6, 0, 4, 4, 6, 4, 3, 3, 4, 1],
  [6, 4, 2, 5, 6, 6, 1, 1, 4, 3, 4, 6, 3],
]
determinant_modulo_7 = 1
\end{verbatim}
\endgroup

\section{Universal norms and one-prime ramification at seven}
\label{seven:universal-norms}

The finite strict module of Theorem~\ref{seven:main} does not coincide
with the Iwasawa module in which ramification is allowed at one of the
two primes above seven. We determine the latter module and its unit
quotient. The distinction comes from capitulation: the faithful class
groups remain of order seven under norms, but their extension maps
vanish, and the resulting defect in norms of global units persists
throughout the tower.

Fix a prime $\mathfrak p\mid7$ of $K=\mathbb Q(\omega)$ and a faithful
character $\chi$ of $\Delta=\operatorname{Gal}(E/K)$. Choose
$\gamma\in\Gamma=\operatorname{Gal}(E_\infty/E)$ with
$\langle\varepsilon_7\rangle(\gamma)=8$, and put
\[
 \Lambda=\mathbb Z_7[[T]],\qquad T=\gamma-1,\qquad
 \omega_n(T)=(1+T)^{7^n}-1.
\]
Let $L_n$ be the completion of $E_n$ at its unique prime over
$\mathfrak p$, and define
\[
 U_n=U^1(L_n),\qquad
 \mathcal E_n=\mathbb Z_7\otimes_{\mathbb Z}\mathcal O_{E_n}^{\times},
 \qquad
 U_\infty=\varprojlim_N U_n,\qquad
 \mathcal E_\infty=\varprojlim_N\mathcal E_n.
\]
The localization map $\ell_n:\mathcal E_n\to U_n$ uses the principal-unit
projection. Its image $\overline{\mathcal E}_n$ is the closure of the
global units which are principal at this prime. Write
$\overline{\mathcal E}_\infty=\varprojlim_N\overline{\mathcal E}_n$ and
\[
 V_{\mathfrak p,\chi}
   =(U_\infty/\overline{\mathcal E}_\infty)^\chi.
\]
Finally, let $X_{\mathfrak p,\infty}$ be the Galois group of the maximal
abelian pro-seven extension of $E_\infty$ unramified outside the primes
over $\mathfrak p$. These are the one-prime local-unit and Galois modules
of \cite[printed p.~36]{Rubin}.

\begin{theorem}\label{seven:one-prime-linear}
For each choice of $\mathfrak p$ and faithful $\chi$, there is
$a_{\mathfrak p,\chi}\in7\mathbb Z_7^\times$ such that
\[
 U_\infty^\chi\simeq\Lambda,\qquad
 \mathcal E_\infty^\chi\simeq\Lambda,\qquad
 V_{\mathfrak p,\chi}\simeq
 X_{\mathfrak p,\infty}^\chi
 \simeq\Lambda/(T-a_{\mathfrak p,\chi}).
\]
The class-field exact sequence is
\begin{equation*}\tag{U1}\label{seven:one-prime-exact}
 0\longrightarrow V_{\mathfrak p,\chi}
 \longrightarrow X_{\mathfrak p,\infty}^\chi
 \longrightarrow M^\chi
 \longrightarrow0,
\end{equation*}
where $M^\chi=\Lambda/(7,T)$, and its first map identifies
$V_{\mathfrak p,\chi}$ with $7X_{\mathfrak p,\infty}^\chi$.
In particular the one-prime characteristic ideal has
$\mu=0$, $\lambda=1$, and is not the unit ideal.

Put
\[
 e_{\mathfrak p,\chi}^{*}
 =-\frac{\log(1+a_{\mathfrak p,\chi})}{\log8}
 \in\mathbb Z_7^\times.
\]
For every generator $f_{\mathfrak p,\chi}$ of
$\operatorname{char}_{\Lambda}X_{\mathfrak p,\infty}^\chi$ and every
$e\in\mathbb Z_7$,
\begin{equation*}\tag{U2}\label{seven:one-prime-weight-valuation}
 v_7\!\left(f_{\mathfrak p,\chi}(8^{-e}-1)\right)
 =1+v_7(e-e_{\mathfrak p,\chi}^{*}),
\end{equation*}
with both sides interpreted as infinity at
$e=e_{\mathfrak p,\chi}^{*}$. These valuations are unbounded on the
positive integer weights, whereas their strict Selmer lengths in
Theorem~\ref{seven:main} are all one.
\end{theorem}

\subsection{Capitulation and norms of global units}

All module notation for completed units in this proof is additive.
Theorem~\ref{seven:main} gives
$A(E_n)^\chi\simeq C_7$ and norm isomorphisms between these groups.
For $n>m$, extension and norm of ideals satisfy
\[
 N_{n/m}j_{n/m}=7^{n-m}
 \quad\hbox{on } A(E_m)^\chi.
\]
Since the right side vanishes and the norm is injective, we have
\begin{equation*}\tag{U3}\label{seven:faithful-capitulation}
 j_{n/m}=0.
\end{equation*}

\begin{lemma}\label{seven:unit-norm-index}
For every $n>m$,
\[
 H^1(\operatorname{Gal}(E_n/E_m),\mathcal E_n)^\chi\simeq C_7,
 \qquad
 [\mathcal E_m^\chi:N_{n/m}\mathcal E_n^\chi]=7.
\]
In particular
$N_{n/0}\mathcal E_n^\chi=7\mathcal E_0^\chi$ for $n\ge1$.
\end{lemma}

\begin{proof}
Set $L=E_n$, $k=E_m$, and $G=\operatorname{Gal}(L/k)$. Hilbert 90 and
the principal-ideal sequence identify
$H^1(G,\mathcal O_L^\times)$ with $P_L^G/P_k$.
The natural map
\[
 P_L^G/P_k\longrightarrow I_L^G/jI_k
 \simeq
 \bigoplus_{\mathfrak q\ {\rm ramified\ in}\ L/k}
           \mathbb Z/e_{\mathfrak q}\mathbb Z
\]
has kernel equal to the capitulation kernel on ideal classes.
The only ramified primes are the two primes above seven, and each is
fixed by $\Delta$. The displayed direct sum is therefore
$\Delta$-trivial. Applying a faithful idempotent and then completing
at seven proves
\[
 H^1(G,\mathcal E_n)^\chi
 \simeq\ker(j_{n/m})^\chi\simeq C_7.
\]
Completion commutes with this cohomology because $\mathbb Z_7$ is flat
over $\mathbb Z$ and the group is finite.

The complex places of $E_n$ form a regular set for
$\Delta\times\operatorname{Gal}(E_n/E)$. Indeed, each complex place has
exactly one embedding extending a fixed embedding of the imaginary
quadratic field $K$. Dirichlet's unit theorem consequently gives the
regular representation with one trivial representation removed.
On a faithful character the removed term disappears, so
\[
 \mathbb Q_7\otimes_{\mathbb Z_7}\mathcal E_n^\chi
 \simeq
 \mathbb Q_7[\operatorname{Gal}(E_n/E)].
\]
Its restriction to $G$ is a sum of regular representations. Hence the
Herbrand quotient of $\mathcal E_n^\chi$ is one: an integral regular
lattice has trivial Tate cohomology, and replacing a lattice by a
commensurable one does not change its Herbrand quotient, since finite
modules have quotient one. Thus
\[
 \#\widehat H^0(G,\mathcal E_n^\chi)
 =\#H^1(G,\mathcal E_n^\chi)=7.
\]
The invariant units are $\mathcal E_m^\chi$, proving the asserted
norm index. At level zero $\mathcal E_0^\chi$ has rank one, so its
submodule of index seven is $7\mathcal E_0^\chi$.
\end{proof}

We record explicitly the base localization. Put
\[
 \varepsilon_0=-270+739\alpha-108\alpha^2.
\]
Reduction using $\alpha^3=269$ gives
\[
 \varepsilon_0^{342}
 \equiv1+7\alpha+28\alpha^2\pmod{49}.
\]
Since $342\equiv6\pmod7$, the principal logarithm therefore satisfies
\begin{equation*}\tag{U4}\label{seven:base-local-generator}
 \frac{\log\varepsilon_0}{7}
 \equiv6\alpha+3\alpha^2\pmod7.
\end{equation*}
The field $L_0$ is the unramified cubic extension of $\mathbb Q_7$,
and logarithm identifies $U_0$ with $7\mathbb Z_7[\alpha]$.
The two faithful lines have bases $7\alpha$ and $7\alpha^2$.
Both coefficients in \eqref{seven:base-local-generator} are units.
As $\mathcal E_0^\chi$ has rank one, this proves
\begin{equation*}\tag{U5}\label{seven:base-local-isomorphism}
 \ell_0:\mathcal E_0^\chi\xrightarrow{\ \sim\ }U_0^\chi.
\end{equation*}

\subsection{Inverse limits and their control maps}

We use two elementary facts about compact modules. First, inverse
limit is exact on inverse systems of compact Hausdorff abelian groups.
To lift a compatible element of a quotient system, any finite set of
compatibility conditions can be met by lifting at its highest level
and taking images at lower levels. Compactness then supplies a lift
meeting all conditions. Transition maps need not be surjective.
In particular, for a compatible endomorphism $T$,
\begin{equation*}\tag{U6}\label{seven:compact-coinvariants}
 (\varprojlim B_n)/T(\varprojlim B_n)
 \simeq\varprojlim(B_n/TB_n).
\end{equation*}
Indeed, exactness applied first to $B_n\to TB_n$ proves
$\varprojlim TB_n=T\varprojlim B_n$, and a second application proves
\eqref{seven:compact-coinvariants}.

Second, the following compact form of Nakayama does not presuppose
finite generation. If a compact pro-seven $\Lambda$-module $B$ has
continuous finite $\Gamma$-stable quotients separating its points and
$B/(7,T)B$ is generated by the image of $b$, then $B=\Lambda b$.
The submodule $\Lambda b$ is compact and closed. Its quotient $Q$
satisfies $Q=(7,T)Q$. Every continuous finite quotient of $Q$ is a
module for some $\mathbb Z_7[C_{7^r}]$ on which $(7,T)$ is nilpotent,
and so must vanish. The finite-level unit groups modulo powers of
seven give the separating quotients needed here.

\begin{lemma}\label{seven:local-unit-control}
For $n>m$, norm induces an isomorphism
\[
 U_n^\chi/(\gamma^{7^m}-1)U_n^\chi
 \xrightarrow{\ \sim\ }U_m^\chi.
\]
Moreover $U_\infty^\chi\simeq\Lambda$, and
$U_\infty^\chi/\omega_mU_\infty^\chi\simeq U_m^\chi$.
\end{lemma}

\begin{proof}
The local extension $L_n/L_m$ is cyclic and totally ramified.
Local class field theory identifies
$U_m/N_{n/m}U_n$ with its Galois group. The $\Delta$-action on that
quotient is trivial because the unramified cubic and cyclotomic
ramified Galois groups commute. Hilbert 90 and the valuation sequence
also give
\[
 H^1(\operatorname{Gal}(L_n/L_m),\mathcal O_{L_n}^\times)
 \simeq\mathbb Z/7^{n-m}\mathbb Z
\]
with trivial $\Delta$-action. Replacing full units by principal units
does not change these Tate groups, since the residue unit group has
order $342$, prime to seven. Both
$\widehat H^0$ and $H^1$ therefore vanish on the faithful component.
The claimed finite norm isomorphism follows from the cyclic Tate
complex.

Compact exactness gives
$U_\infty^\chi/TU_\infty^\chi\simeq U_0^\chi\simeq\mathbb Z_7$.
Compact Nakayama makes $U_\infty^\chi$ cyclic over $\Lambda$.
All the finite local norms on this character are surjective, so a
generator of the limit maps to a generator of every $U_n^\chi$.
These groups are free of rank $7^n$ over $\mathbb Z_7$. There are no
seven-power roots of unity, since the local ramification index is a
power of seven, and their rational rank follows from the local
normal-basis representation and logarithm. Consequently the
surjection
\[
 \Lambda/\omega_n\Lambda\longrightarrow U_n^\chi
\]
is an isomorphism between free $\mathbb Z_7$-modules of the same rank.
The annihilator of a generator of $U_\infty^\chi$ lies in every
$\omega_n\Lambda$. But $\omega_n\in(7,T)^{n+1}$, and their intersection
is zero. This proves freeness. The same finite norm argument and
compact exactness give control at every level $m$.
\end{proof}

The global control map requires the capitulation terms, including
their transition maps.
For $G_n=\operatorname{Gal}(E_n/E)$ and $n\ge1$, the cyclic Tate
complex and Lemma~\ref{seven:unit-norm-index} give
\begin{equation*}\tag{U7}\label{seven:global-coinvariant-sequence}
 0\longrightarrow H^1(G_n,\mathcal E_n)^\chi
 \longrightarrow\mathcal E_n^\chi/T\mathcal E_n^\chi
 \xrightarrow{\ N_{n/0}\ }7\mathcal E_0^\chi
 \longrightarrow0.
\end{equation*}
The left term has order seven, but its norm transition from any higher
positive level to a lower one is zero. To see this, represent a base
class by an ideal $\mathfrak a$ and choose
$(z_n)=\mathfrak a\mathcal O_{E_n}$. Its cohomology class is represented
by $\gamma(z_n)/z_n$. For $n>m$, put $z_m'=N_{n/m}z_n$.
Then
\[
 (z_m')=\mathfrak a^{7^{n-m}}\mathcal O_{E_m}.
\]
The base class has order dividing seven, so
$\mathfrak a^{7^{n-m}}=(b)$ for some $b\in E^\times$.
Thus $z_m'/b$ is a unit, and
\[
 N_{n/m}\!\left(\frac{\gamma(z_n)}{z_n}\right)
 =\frac{\gamma(z_m'/b)}{z_m'/b}
\]
is a coboundary. This calculation commutes with the faithful
idempotent. The transition on the right of
\eqref{seven:global-coinvariant-sequence} is the identity.
Taking compact inverse limits therefore proves
\begin{equation*}\tag{U8}\label{seven:global-limit-coinvariants}
 \mathcal E_\infty^\chi/T\mathcal E_\infty^\chi
 \xrightarrow{\ \sim\ }7\mathcal E_0^\chi.
\end{equation*}
In particular there is no unaccounted kernel in this specialization.
Compact Nakayama now proves that $\mathcal E_\infty^\chi$ is cyclic.

The finite maps $\mathcal E_n\to\overline{\mathcal E}_n$ are surjective
maps of compact groups and commute with norms. Compact exactness
identifies $\overline{\mathcal E}_\infty$ with the image of
$\mathcal E_\infty\to U_\infty$, without requiring injectivity of
finite-level localization. Choose cyclic generators $e_\infty$ and
$u_\infty$ of $\mathcal E_\infty^\chi$ and $U_\infty^\chi$, and write
\[
 \ell_\infty(e_\infty)=f_{\mathfrak p,\chi}(T)u_\infty.
\]
Equations \eqref{seven:base-local-isomorphism} and
\eqref{seven:global-limit-coinvariants} give
\begin{equation*}\tag{U9}\label{seven:unit-factor-augmentation}
 f_{\mathfrak p,\chi}(0)\in7\mathbb Z_7^\times.
\end{equation*}
Thus $f_{\mathfrak p,\chi}\ne0$. Any relation on $e_\infty$ would
annihilate this nonzero element in the free local module $\Lambda$.
As $\Lambda$ is a domain, the global cyclic module is free and its
localization is injective. We have proved
\begin{equation*}\tag{U10}\label{seven:unit-factor}
 \mathcal E_\infty^\chi\simeq\Lambda,\qquad
 V_{\mathfrak p,\chi}\simeq\Lambda/(f_{\mathfrak p,\chi}).
\end{equation*}
The augmentation condition alone allows either $\mu=1,\lambda=0$
or $\mu=0,\lambda\ge1$. A first-layer calculation distinguishes them.

\subsection{An exact first-layer local index}

Use the period $\beta$, the automorphism $\tau$, and the thirteen
units of \S\ref{seven:finite-data}. Put $\pi=\beta+1$. The period
polynomial gives the Eisenstein equation
\begin{equation*}\tag{U11}\label{seven:local-eisenstein}
 \pi^7-7\pi^6+49\pi^4-98\pi^2-49\pi+7=0.
\end{equation*}
Thus $v_\pi(7)=7$, the residue field is
$\mathbb F_7[\overline\alpha]$ with $\overline\alpha^3=3$, and
$7/\pi^7\equiv-1\pmod\pi$. Set $U^j=1+\pi^j\mathcal O_{L_1}$.

\begin{lemma}\label{seven:local-finite-basis}
Let $\chi_a(\delta)=\zeta_3^a$, where
$\delta(\alpha)=\zeta_3\alpha$, $a=1,2$, and
$\zeta_3\equiv2\pmod7$. For
$J=\{1,2,3,4,5,6,8\}$, the seven elements
\[
 b_{a,j}=e_{\chi_a}(1+\alpha^a\pi^j),\qquad j\in J,
\]
form a $\mathbb Z_7$-basis of $U_1^{\chi_a}$.
Their images in $(U^1/U^{16})^{\chi_a}$ have respective orders
$343,49,49,49,49,49,49$ and form a direct-product basis.
\end{lemma}

\begin{proof}
For $i=0,1,2$, the successive leading grades of
$g_i=1+\alpha^i\pi$, $g_i^7$, and $g_i^{49}$ are $1,7,14$,
with coefficients
$\overline\alpha^i,\overline\alpha^{7i},-\overline\alpha^{7i}$.
For $h_{i,j}=1+\alpha^i\pi^j$, with
$j\in\{2,3,4,5,6,8\}$, the leading grades of $h_{i,j}$ and
$h_{i,j}^7$ are $j,j+7$, with coefficients
$\overline\alpha^i,-\overline\alpha^i$.
These assertions follow from the binomial expansion: the competing
lowest grades in $(1+c\pi^j)^7-1$ are $7+j$ and $7j$.
The next powers lie in $U^{16}$.

The displayed pivot grades partition $\{1,\ldots,15\}$, and the three
coefficients at each grade form an $\mathbb F_7$-basis of the residue
field. Successive digit elimination gives
\[
 U^1/U^{16}\simeq(C_{343})^3\oplus(C_{49})^{18}
\]
with these generators. On each associated graded quotient,
$e_{\chi_a}$ projects onto
$\mathbb F_7\overline\alpha^a$.
Moreover
$\overline\alpha^{7a}=3^{2a}\overline\alpha^a$ stays in that line.
The same elimination proves the asserted faithful direct-product
basis.

Logarithm and exponential are inverse on $U^j$ for $j\ge2$, since
$j>7/(7-1)$. Hence
$U^{16}=(U^9)^7\subseteq(U^1)^7$.
Reduction modulo $U^{16}$ therefore preserves
$U_1^{\chi_a}/7U_1^{\chi_a}$. The seven elements generate this quotient,
and $U_1^{\chi_a}$ is free of rank seven over $\mathbb Z_7$.
Nakayama proves that the displayed lifts form a full basis.
\end{proof}

The seven projected global units to compare with this basis are
\[
 e_{\chi_a}\varepsilon_0,\quad
 e_{\chi_a}\eta,\quad e_{\chi_a}\tau\eta,\quad\ldots,\quad
 e_{\chi_a}\tau^5\eta.
\]
They form a full $\mathbb Z_7$-basis of $\mathcal E_1^{\chi_a}$.
Indeed, prime-to-seven descent identifies
$\mathcal E(F_1)=\mathcal E(E_1)^{c=1}$, and projection gives an exact
sequence
\[
 0\longrightarrow\mathcal E(Q_1)
 \longrightarrow\mathcal E(F_1)
 \xrightarrow{\ e_{\chi_a}\ }\mathcal E(E_1)^{\chi_a}
 \longrightarrow0.
\]
For the kernel, the $\Delta$-trivial unit module is
$\mathcal E(KQ_1)$; its units are fixed by $c$ after seven-adic
completion. In fact $u/c(u)$ has absolute value one at all complex
embeddings and is a root of unity, and norm and extension identify
the completed units of this CM field with those of $Q_1$.
For surjectivity a vector $y$ in the faithful line lifts to $y+cy$.
The first six units in the thirteen-element seven-saturated basis
lie in the rank-six kernel. Their span is a direct summand, and hence
saturated, so it is the entire kernel. The remaining seven vectors
therefore give the stated basis.

Raise these seven global units to $342$ to make them principal.
This multiplies their completed coordinate columns by a
$\mathbb Z_7$-unit, and does not change their lattice index.
In the local basis of Lemma~\ref{seven:local-finite-basis}, the
coordinate matrices modulo $49$, with row order
$1,2,3,4,5,6,8$ and column order
$\varepsilon_0,\eta,\tau\eta,\ldots,\tau^5\eta$, are
\begin{equation*}\tag{U12}\label{seven:local-matrix-one}
 B_1=
 \begin{pmatrix}
 21&0&28&14&35&14&35\\
 0&48&48&34&48&27&48\\
 0&17&12&0&2&11&27\\
 28&45&24&25&41&2&34\\
 0&48&39&34&29&48&3\\
 7&26&15&30&16&22&12\\
 3&44&4&40&9&8&0
 \end{pmatrix},
 \qquad\det B_1\equiv28\pmod{49},
\end{equation*}
and
\begin{equation*}\tag{U13}\label{seven:local-matrix-two}
 B_2=
 \begin{pmatrix}
 42&42&21&0&14&21&28\\
 0&30&23&9&9&16&37\\
 0&20&44&33&29&32&28\\
 21&37&43&19&42&14&33\\
 0&22&32&31&13&14&28\\
 28&11&29&0&35&23&37\\
 13&44&0&12&34&23&32
 \end{pmatrix},
 \qquad\det B_2\equiv7\pmod{49}.
\end{equation*}

For precision, the finite arithmetic needed to verify these entries
can be performed entirely in successive $\pi$-adic digits. Equation
\eqref{seven:local-eisenstein} gives the exact carry rule
\begin{equation*}\tag{U14}\label{seven:local-carry}
 7\equiv-\pi^7-\pi^{13}+\pi^{15}\pmod{\pi^{16}}.
\end{equation*}
Use sixteen digits, each in the basis
$1,\alpha,\alpha^2$ with coefficients in $\{0,\ldots,6\}$.
Reduce $\alpha^3=269$ and carry every coefficient multiple of seven
upward by \eqref{seven:local-carry}; each carry strictly increases its
$\pi$-degree. This gives the ring
$\mathcal O_{L_1}/(\pi^{16})$ of order $7^{48}$.
All denominators in the printed unit expressions are prime to seven,
and their inverses modulo $343$ suffice, since $7^3=0$ in this ring.
Use $\beta=\pi-1$, the printed polynomial for $\tau$, and
$\zeta_3=324\pmod{343}$. On a principal unit the projector is the
product of its three $\Delta$-conjugates with exponents
$3^{-1}\zeta_3^{-ar}\pmod{343}$, $r=0,1,2$.
The pivot grades in the preceding lemma then recover the first
coordinate modulo $343$ and the other six modulo $49$.
This procedure determines all entries of
\eqref{seven:local-matrix-one}--\eqref{seven:local-matrix-two} from the
unit data already printed in \S\ref{seven:finite-data}.

Both determinants have valuation exactly one. The local and global
vectors used are full $\mathbb Z_7$-bases, so we have proved the exact
indices
\begin{equation*}\tag{U15}\label{seven:first-layer-local-index}
 [U_1^\chi:\overline{\mathcal E}_1^\chi]=7
\end{equation*}
for both faithful characters. The calculation applies at either
chosen prime above seven; changing the image of $\omega$ interchanges
the character labels.
As a further exact norm check on the same data, the printed expression
$\eta(\alpha,\beta)$ satisfies
\[
 \operatorname{Res}_B(q(B),\eta(\alpha,B))=1
 \quad\hbox{in }\mathbb Q[\alpha]/(\alpha^3-269).
\]
The six $Q_1$-units have relative norms
$-1,-1,1,-1,-1,1$, and $\varepsilon_0$ has relative norm
$\varepsilon_0^7$. This confirms the first-layer global norm defect
without serving as an additional premise for the all-level argument.

\subsection{The linear factor and the weight valuations}

For every $n>1$, Lemma~\ref{seven:unit-norm-index} gives
$[\mathcal E_1^\chi:N_{n/1}\mathcal E_n^\chi]=7$.
These norm images form a descending chain with the same index, hence
are equal. Compactness identifies the common image with the
projection of $\mathcal E_\infty^\chi$ to level one.
Lemma~\ref{seven:local-unit-control} therefore gives
\[
 V_{\mathfrak p,\chi}/\omega_1V_{\mathfrak p,\chi}
 \simeq
 U_1^\chi/
 \ell_1\bigl(\operatorname{im}
       (\mathcal E_\infty^\chi\to\mathcal E_1^\chi)\bigr).
\]
Equation \eqref{seven:first-layer-local-index} also proves that
$\ell_1$ is injective on this character, since its source and target
have equal rank. The universal-norm image has index seven in the
global units, so
\begin{equation*}\tag{U16}\label{seven:unit-factor-first-quotient}
 \#\Lambda/(f_{\mathfrak p,\chi},\omega_1)=7^2.
\end{equation*}

We now determine the Weierstrass invariants. If all coefficients of
$f=f_{\mathfrak p,\chi}$ are divisible by seven, then
\eqref{seven:unit-factor-augmentation} implies $f=7u$ for a unit $u$.
The quotient in \eqref{seven:unit-factor-first-quotient} would then
have order $7^7$, a contradiction. Otherwise Weierstrass preparation
writes $f=ug$, where $g$ is a distinguished polynomial of degree
$d\ge1$. Its constant coefficient has valuation one, so $g$ is
Eisenstein. If $1\le d\le5$, then at a primitive seventh root
$\zeta$ the leading term of $g(\zeta-1)$ has uniquely smallest
valuation $d/6$. Including the root $\zeta=1$, we obtain
\[
 v_7\!\left(\prod_{\zeta^7=1}g(\zeta-1)\right)=1+d.
\]
If $d\ge6$, each nontrivial-root factor has valuation at least one,
so the displayed valuation is at least seven, or the quotient is
infinite. Equation \eqref{seven:unit-factor-first-quotient} therefore
forces $d=1$. We have proved
\[
 f_{\mathfrak p,\chi}(T)
 =u(T)(T-a_{\mathfrak p,\chi}),\qquad
 a_{\mathfrak p,\chi}\in7\mathbb Z_7^\times.
\]

Rubin's class-field sequence
\cite[printed p.~36, equation~(5)]{Rubin} gives
\eqref{seven:one-prime-exact}. Its quotient is the finite faithful
class module of Theorem~\ref{seven:main}. Thus
$X_{\mathfrak p,\infty}^\chi$ is a finitely generated
$\mathbb Z_7$-module of rank one.
By \cite[Theorem~5.3(v), printed pp.~40--41]{Rubin},
$X_{\mathfrak p,\infty}$ has no nonzero pseudo-null submodule;
in the one-variable setting this excludes every finite submodule.
The $\mathbb Z_7$-torsion of $X_{\mathfrak p,\infty}^\chi$ is finite
and $\Lambda$-stable, so it vanishes. Hence this module is
$\mathbb Z_7$-free of rank one. Its $\Gamma$-action agrees, on the
finite-index submodule $V_{\mathfrak p,\chi}$, with multiplication
by $1+a_{\mathfrak p,\chi}$, and must agree on the whole module.
The submodule of index seven is consequently
$7X_{\mathfrak p,\infty}^\chi$. This proves all the asserted module
identifications.

The hypotheses of Rubin's theorem are satisfied: $K$ has class
number one and six roots of unity, $E/K$ is abelian of degree three,
and $E_\infty/K$ is abelian with
$\operatorname{Gal}(E_\infty/E)\simeq\mathbb Z_7$.
The statement concerns any such one-variable extension, so it
applies to the cyclotomic line; it is not restricted to the
$\mathfrak p$-ramified line.

It remains to evaluate the linear factor. Both
$\log(1+a_{\mathfrak p,\chi})$ and $\log8$ have valuation one, so
$e_{\mathfrak p,\chi}^{*}$ is a seven-adic unit and
$8^{-e_{\mathfrak p,\chi}^{*}}=1+a_{\mathfrak p,\chi}$.
For every $e\in\mathbb Z_7$,
\[
 8^{-e}-(1+a_{\mathfrak p,\chi})
 =(1+a_{\mathfrak p,\chi})
   \left(8^{-(e-e_{\mathfrak p,\chi}^{*})}-1\right).
\]
The logarithm and exponential on $7\mathbb Z_7$ give
$v_7(8^z-1)=1+v_7(z)$ for $z\in\mathbb Z_7$.
A unit in $\Lambda$ evaluates to a unit at every point of
$7\mathbb Z_7$. Equation
\eqref{seven:one-prime-weight-valuation} follows for every
characteristic generator.
Positive integers approximate any element of $\mathbb Z_7$ to
arbitrary precision, and hence these valuations are unbounded.

The standard Iwasawa involution replaces the root
$a_{\mathfrak p,\chi}$ by
$(1+a_{\mathfrak p,\chi})^{-1}-1$, which still has valuation one,
and replaces $e_{\mathfrak p,\chi}^{*}$ by its negative.
It does not remove the discrepancy with the constant strict length.
Thus identifying the strict length with the valuation of a generator
of this one-prime characteristic ideal is false, independently of
unit ambiguity and this involution.
An identification with a particular normalized Katz function still
requires its precise character, period, reciprocity, and Euler-factor
conventions. The conclusion here concerns the characteristic
generator of the explicitly defined one-prime module.


\section*{Method and availability}
The text and programs were produced by GPT-6 Astra in ChatGPT Pro mode.
The analytic and arithmetic programs are included in the project.
The seven-primary appendices contain their complete proofs and exact witnesses, including two independent exact calculations of the first-layer local-unit matrices.
An independently editable extract includes the same proof and portable verifiers.
The mathematical checks include a statement-only partial reconstruction, a second statement-only reconstruction of the full argument with the four finite arithmetic assertions granted, source-based audits, separate same-model checks of the specialization and ray arguments, and fresh computations proving those arithmetic assertions; their distinct scopes and outcomes are recorded in the reproducibility material accompanying this paper.
The workbench address is \expandafter\url\expandafter{\workbenchurl}, and the project record address is \expandafter\url\expandafter{\recordurl}.
No human verification is claimed.
The original text and programs are dedicated to the public domain under CC0~1.0 Universal; cited publications retain their respective rights.

\begin{thebibliography}{99}
\bibitem[Belabas and Jaulent(2016)]{BelabasJaulent2016}
Belabas, K., \& Jaulent, J.-F. (2016). The logarithmic class group package in PARI/GP.
\emph{Publications math\'ematiques de Besan\c con. Alg\`ebre et th\'eorie des nombres}, 5--18.
\url{https://doi.org/10.5802/pmb.o-1}.
Author version: \url{https://www.math.u-bordeaux.fr/~kbelabas/research/bnflog-2016.pdf}.
\bibitem[Dunn and Radziwi\l\l(2024)]{DunnRadziwill2024}
Dunn, A., \& Radziwi\l\l, M. (2024). Bias in cubic Gauss sums: Patterson's conjecture. \emph{Annals of Mathematics, 200}(3), 967--1057. \url{https://doi.org/10.4007/annals.2024.200.3.3}. Theorem~1.1; cited locator from \href{https://arxiv.org/abs/2109.07463v3}{arXiv version~3}.

\bibitem[Jaulent and Michel(2006)]{JaulentMichel2006}
Jaulent, J.-F., \& Michel, A. (2006). Approche logarithmique des noyaux \'etales sauvages des corps de nombres.
\emph{Journal of Number Theory, 120}(1), 72--91.
\url{https://www.math.u-bordeaux.fr/~jjaulent/Articles/J44.pdf}.
\bibitem[Kolster and Movahhedi(2000)]{KolsterMovahhedi2000}
Kolster, M., \& Movahhedi, A. (2000). Galois co-descent for \'etale wild kernels and capitulation.
\emph{Annales de l'Institut Fourier, 50}(1), 35--65.
\url{https://doi.org/10.5802/aif.1746}.
\bibitem[Louboutin and Munsch(2013)]{LouboutinMunsch2013}
Louboutin, S. R., \& Munsch, M. (2013).
On positive real zeros of theta and $L$-functions associated with real, even and primitive characters.
\emph{Publicationes Mathematicae Debrecen, 83}(4), 643--665.
Theorem 4, pp.~646--647.
\url{https://publi.math.unideb.hu/paper/1837/download/10_5486_PMD_2013_5670.pdf}.
\bibitem[Milne(2020a)]{MilneANT2020}
Milne, J. S. (2020a). \emph{Algebraic number theory} (Version 3.08).
\url{https://www.jmilne.org/math/CourseNotes/ANT.pdf}.
\bibitem[Milne(2020b)]{Milne2020}
Milne, J. S. (2020b). \emph{Class field theory} (Version 4.03).
\par\noindent\url{https://www.jmilne.org/math/CourseNotes/CFT.pdf}.
\bibitem[OpenAI(2026a)]{OpenAIFirst2026}
OpenAI. (2026a, September 25). \emph{An unconditional first moment for cubic Gauss sums}. Theorem~1.1, p.~1. \href{https://github.com/openai/math/blob/main/preprints/An-unconditional-first-moment-for-cubic-Gauss-sums-September-25-2026/paper.pdf}{Dated primary release}.

\bibitem[OpenAI(2026b)]{OpenAIQuasi2026}
OpenAI. (2026b, September 30). \emph{The quasi-Riemann hypothesis: A zero-free half-plane $\Re s>7/8$}. Theorem~1.1, p.~4. \href{https://github.com/openai/math/blob/main/preprints/The-Quasi-Riemann-Hypothesis-September-30-2026/paper.pdf}{Dated primary release}.

\bibitem[Poonen(2015)]{Poonen2015}
Poonen, B. (2015). \emph{Tate's thesis} (MIT 18.786 lecture notes).
\url{https://math.mit.edu/~poonen/786/notes.pdf}.
\bibitem[Rubin(1991)]{Rubin}
Rubin, K. (1991). The ``main conjectures'' of Iwasawa theory for imaginary quadratic fields.
\emph{Inventiones Mathematicae, 103}, 25--68.
\url{https://doi.org/10.1007/BF01239508}.
\bibitem[Schneider(1979)]{Schneider1979}
Schneider, P. (1979). \"Uber gewisse Galoiscohomologiegruppen.
\emph{Mathematische Zeitschrift, 168}, 181--205.
\url{https://gdz.sub.uni-goettingen.de/id/PPN266833020_0168}.
\bibitem[Weibel(2013)]{Weibel2013}
Weibel, C. A. (2013). \emph{The $K$-book: An introduction to algebraic $K$-theory}.
Graduate Studies in Mathematics, 145. American Mathematical Society.
\url{https://sites.math.rutgers.edu/~weibel/Kbook/Kbook.VI.pdf}.
\bibitem[Yu(2014)]{Yu2014}
Yu, C.-F. (2014). \emph{Ambiguous class number formulas} (MPIM Preprint 2014-68).
Max Planck Institute for Mathematics.
\url{https://archive.mpim-bonn.mpg.de/id/eprint/3296/1/preprint_2014_68.pdf}.
\bibitem[Fukuda(1994)]{sevenFukuda1994}
Fukuda, T. (1994).
Remarks on $\mathbb Z_p$-extensions of number fields.
\emph{Proceedings of the Japan Academy, Series A, Mathematical Sciences},
\emph{70}(8), 264--266.
\url{https://doi.org/10.3792/pjaa.70.264}.

\bibitem[Friedman(1989)]{sevenFriedman1989}
Friedman, E. (1989).
Analytic formulas for the regulator of a number field.
\emph{Inventiones Mathematicae}, \emph{98}(3), 599--622.
\url{https://doi.org/10.1007/BF01393839}.

\bibitem[Sutherland(2015)]{sevenSutherland2015}
Sutherland, A. V. (2015, November 12).
\emph{The analytic class number formula} (18.785 Lecture Notes 18).
Massachusetts Institute of Technology.
\url{https://math.mit.edu/classes/18.785/2015fa/LectureNotes18.pdf}.

\bibitem[Simard(2015)]{sevenSimard2015}
Simard, N. (2015, November 2).
\emph{Notes on theta series}.
\url{https://www.math.mcgill.ca/goren/GS/Thetas.pdf}.

\bibitem[Darmon et al.(2015)]{sevenDarmonLauderRotger2015}
Darmon, H., Lauder, A., \& Rotger, V. (2015).
Stark points and $p$-adic iterated integrals attached to modular forms of
weight one.
\emph{Forum of Mathematics, Pi}, \emph{3}, e8.
\url{https://doi.org/10.1017/fmp.2015.7}.

\bibitem[Johansson(2017)]{sevenJohansson2017}
Johansson, F. (2017).
Arb: Efficient arbitrary-precision midpoint-radius interval arithmetic.
\emph{IEEE Transactions on Computers}, \emph{66}(8), 1281--1292.
\url{https://doi.org/10.1109/TC.2017.2690633}.

\bibitem[Sharifi(n.d.)]{sevenSharifi}
Sharifi, R. (n.d.).
\emph{Iwasawa theory} [Author's lecture notes].
\url{https://www.math.ucla.edu/~sharifi/notes/iwasawa.pdf}.
Chapters 1 and 3, especially Corollaries 1.3.9 and 1.3.12,
Remark 1.3.21, and Proposition 3.2.20.

\end{thebibliography}
\end{document}
